%%%%%%%%%%%%%%%%%%%%%%%%%%%%%%%%%%%%%%%%%%%%%%%%%%%%%%%%%%%%%%%%%%%%%%%%%%%%%%%
% CHAPTER: K3 x E AND THE CY3 CHIRAL PROGRAMME
% Split from toroidal_elliptic.tex (lines 1496-4807)
% Contains: CY gluing, K3 chiral algebra, K3xE threefold, BPS, BKM, 
%           relative chiral, ADE Yangian, anomaly, M-theory web
%%%%%%%%%%%%%%%%%%%%%%%%%%%%%%%%%%%%%%%%%%%%%%%%%%%%%%%%%%%%%%%%%%%%%%%%%%%%%%%

%% ===============================================================
%% CALABI--YAU LOCAL-TO-GLOBAL GLUING AND THE K3 CHIRAL ALGEBRA
%% ===============================================================

\section{Calabi--Yau local-to-global gluing}
\label{sec:cy-local-global}
\index{Calabi--Yau!local-to-global gluing|textbf}
\index{McKay correspondence}
\index{Ginzburg dg algebra}

A compact Calabi--Yau manifold is assembled from local charts, and the chiral algebra must be assembled the same way. The fundamental tension: locally, each chart of $K3 \times E$ looks like $\C^2/\Gamma \times \C$ for a finite subgroup $\Gamma \subset \mathrm{SL}(2,\C)$, and the local associative algebra is the CoHA of the McKay quiver, a completely explicit object. Globally, these local algebras must be glued by $A_\infty$-bimodule transition data along the overlaps, subject to cocycle conditions on triple overlaps. The quiver-chart description is necessary because no single global construction produces $A_{K3 \times E}$ directly; the CY-to-chiral correspondence $\Phi_3$ of Theorem~\ref{thm:cy-to-chiral-d3} acts on the derived category as a whole, while the explicit computation requires the local description. The resolution requires two ingredients: the local charts (McKay correspondence, Ginzburg dg algebras, preprojective algebras) and the gluing mechanism ($A_\infty$-bimodules, \v{C}ech descent) that assembles them into a global object.

\subsection{McKay correspondence and ADE quiver charts}
\label{subsec:mckay-quiver}
\label{sec:ade-yangian-level1}

A K3 surface $S$ acquires ADE singularities at special points in
moduli. Near each singularity, the local geometry is
$\bC^2/\Gamma$ for a finite subgroup $\Gamma \subset \mathrm{SL}(2,\bC)$.
The McKay correspondence (McKay 1980, Gonzalez-Sprinberg--Verdier
1983) gives a bijection
\[
\bigl\{\text{nontrivial irreps of } \Gamma\bigr\}
\;\xleftrightarrow{\;\sim\;}\;
\bigl\{\text{exceptional curves in the minimal resolution}\bigr\},
\]
and the McKay quiver $Q_\Gamma$ encodes the tensor product rule
$\rho_{\mathrm{fund}} \otimes \rho_i = \bigoplus_j a_{ij}\,\rho_j$,
where $a_{ij}$ are the arrows of the extended ADE Dynkin diagram.

\begin{definition}[McKay quiver and Cartan matrix]
\label{def:mckay-quiver}
\index{McKay quiver|textbf}
For a finite subgroup $\Gamma \subset \mathrm{SL}(2,\bC)$ of ADE
type, the \emph{McKay quiver} $Q_\Gamma$ has vertices indexed by
irreducible representations $\rho_0 = \mathrm{triv},\,
\rho_1,\ldots,\rho_r$ and $a_{ij}$ arrows from vertex $i$ to
vertex $j$, where $a_{ij}$ is the multiplicity of $\rho_j$ in
$\rho_{\mathrm{fund}} \otimes \rho_i$. The extended Cartan matrix is
\[
C_{\mathrm{ext}} = 2I - A,
\]
where $A$ is the adjacency matrix of $Q_\Gamma$.
\end{definition}

\begin{example}[ADE classification]
\label{ex:ade-classification}
\index{ADE classification!finite subgroups of SL(2)}
The finite subgroups of $\mathrm{SL}(2,\bC)$ are classified by
ADE Dynkin diagrams:
\begin{center}
\small
\renewcommand{\arraystretch}{1.2}
\begin{tabular}{lllrl}
\textbf{Group} & $\Gamma$ & \textbf{Type} & $|\Gamma|$
 & \textbf{Irreps (rank)} \\
\hline
Cyclic & $\bZ/n$ & $A_{n-1}$ & $n$ & $n$ irreps, all $1$-dim \\
Binary dihedral & $\mathrm{BD}_n$ & $D_{n+2}$ & $4n$ & $n+3$ irreps \\
Binary tetrahedral & $\mathrm{BT}$ & $E_6$ & $24$ & $7$ irreps \\
Binary octahedral & $\mathrm{BO}$ & $E_7$ & $48$ & $8$ irreps \\
Binary icosahedral & $\mathrm{BI}$ & $E_8$ & $120$ & $9$ irreps
\end{tabular}
\end{center}
For each type, $|\Gamma| = \sum d_i^2$ (sum of squared dimensions
of irreducible representations) provides an independent check of
the character table.
\end{example}

\begin{proposition}[Preprojective algebra from the McKay quiver;
\ClaimStatusProvedElsewhere]
\label{prop:preprojective-mckay}
\index{preprojective algebra|textbf}
The preprojective algebra $\Pi(Q_\Gamma) = kQ_{\mathrm{double}}/
(\sum [a, a^*])$, where $Q_{\mathrm{double}}$ is the double quiver
(adjoin a reverse arrow $a^*$ for each arrow $a$), has centre
isomorphic to the invariant ring:
\[
Z(\Pi(Q_\Gamma)) \;\cong\; \bC[x,y,z]^\Gamma.
\]
For the Dynkin quiver of type~$\Delta$ with $r$ vertices and
Coxeter number~$h$, the preprojective algebra is
finite-dimensional with $\dim \Pi_0(Q_\Delta)$ given by the
following table \textup{(}Crawley-Boevey 2001, computed via
$\dim e_i\Pi e_j = \min(i,j)\min(r+1-i,r+1-j)$ for type~$A$
and by the McKay quiver formula
$\dim \Pi(Q_\Gamma) = |\Gamma| \cdot (r+1)$ for the extended
quiver\textup{)}:
\begin{center}
\small
\renewcommand{\arraystretch}{1.2}
\begin{tabular}{lrrrrr}
\textbf{Type} & $|\Gamma|$ & $r = \operatorname{rank}$
 & $h$ & $|Q_0|$ & $\dim\Pi_0$ \\
\hline
$A_1$ & $2$ & $1$ & $2$ & $2$ & $4$ \\
$A_2$ & $3$ & $2$ & $3$ & $3$ & $9$ \\
$A_3$ & $4$ & $3$ & $4$ & $4$ & $20$ \\
$A_4$ & $5$ & $4$ & $5$ & $5$ & $35$ \\
$D_4$ & $8$ & $4$ & $6$ & $5$ & $40$ \\
$D_5$ & $12$ & $5$ & $8$ & $6$ & $72$ \\
$E_6$ & $24$ & $6$ & $12$ & $7$ & $168$ \\
$E_7$ & $48$ & $7$ & $18$ & $8$ & $384$ \\
$E_8$ & $120$ & $8$ & $30$ & $9$ & $1080$
\end{tabular}
\end{center}
The table entries are computed from the McKay quiver formula
$\dim e_i\Pi e_j = \min(i,j)\min(r{+}1{-}i,\,r{+}1{-}j)$
for type~$A$ and from representation-theoretic data for the
exceptional types.
\end{proposition}

\subsection{Ginzburg dg algebras and Jacobian algebras}
\label{subsec:ginzburg-dg}

\begin{definition}[Ginzburg dg algebra]
\label{def:ginzburg-dg}
\index{Ginzburg dg algebra|textbf}
\index{Jacobian algebra}
For an ADE quiver $Q$ with potential $W$ (a linear combination
of cyclic paths), the \emph{Ginzburg dg algebra} $\cA_\Gamma$
is a Calabi--Yau $3$-fold dg algebra concentrated in degrees
$[-3, 0]$. Its zeroth cohomology is the \emph{Jacobian algebra}:
\[
\mathrm{Jac}(Q, W) = kQ \big/ \bigl(\partial_a W : a \in Q_1\bigr),
\]
where $\partial_a W$ is the cyclic derivative with respect to
arrow $a$. For ADE types, $\mathrm{Jac}(Q,W)$ is isomorphic
to the preprojective algebra $\Pi_0(Q)$.
\end{definition}

\begin{remark}[CY3 extension]
\label{rem:cy3-extension}
For the local CY3 model $\bC^2/\Gamma \times \bC$, one adds a
loop at each vertex of the McKay quiver (from the $\bC$
direction). The extended quiver with potential $(Q_{\mathrm{ext}},
W_{\mathrm{ext}})$ has Jacobian algebra that is $3$-Calabi--Yau
in the sense of Ginzburg: the derived category
$D^b(\mathrm{Jac}(Q_{\mathrm{ext}}, W_{\mathrm{ext}}))$
carries a Serre functor $S = [3]$.
\end{remark}

\subsection{$A_\infty$-bimodule transition data}
\label{subsec:ainfty-bimodule}

\begin{construction}[$A_\infty$-bimodule gluing]
\label{constr:ainfty-bimodule}
\index{A-infinity bimodule!chart transition}
\index{reduced bar complex!bimodule}
Given two local charts $U_\alpha, U_\beta$ with tilting
generators $T_\alpha, T_\beta$ and endomorphism algebras
$A_\alpha = \mathrm{End}(T_\alpha)$,
$A_\beta = \mathrm{End}(T_\beta)$, the transition data on the
overlap $U_{\alpha\beta} = U_\alpha \cap U_\beta$ is encoded
in the $(A_\alpha, A_\beta)$-bimodule
\[
M_{\alpha\beta} = R\mathrm{Hom}(T_\alpha|_{U_{\alpha\beta}},\,
T_\beta|_{U_{\alpha\beta}})
\]
with $A_\infty$-bimodule operations
$m_{p|1|q}\colon A_\alpha^{\otimes p} \otimes M_{\alpha\beta}
\otimes A_\beta^{\otimes q} \to M_{\alpha\beta}[2-p-q]$
satisfying the $A_\infty$-bimodule equation. The derived tensor
product $M_{\alpha\beta} \otimes^{\mathbf{L}}_{A_\beta}
M_{\beta\gamma}$ provides the triple-overlap cocycle data.
The reduced bar complex $\overline{B}(A_\alpha, M_{\alpha\beta},
A_\beta)$ computes the relevant $\mathrm{Ext}$ groups.
\end{construction}

\begin{remark}[Reduced bar complex and unit elements]
\label{rem:reduced-bar-bimodule}
\index{reduced bar complex!augmentation ideal}
\index{A-infinity bimodule!augmentation ideal}
The $A_\infty$-bimodule relations hold on the \emph{reduced}
bar complex (augmentation ideal only): unit elements do not
participate in the higher operations $m_{p|1|q}$ for $p + q \geq 2$.
This is a general structural principle, not specific to K3.
When computing $\overline{B}(A_\alpha, M_{\alpha\beta}, A_\beta)$,
one restricts to the augmentation ideal
$\overline{A} = \ker(\varepsilon \colon A \to \Bbbk)$
in both algebra arguments. The unit $1_A \in A$ acts on~$M$
via $m_{1|1|0}(1_A, m) = m$ (identity bimodule action) and
$m_{0|1|1}(m, 1_B) = m$, but does not enter the higher bar
differential. For the K3 quiver charts, where
$A_\alpha = \mathrm{End}(T_\alpha)$ is augmented over~$\Bbbk$,
this restriction ensures that only the interacting OPE data
contributes to the gluing obstructions.
\end{remark}

\subsection{\v{C}ech descent and the global category}
\label{subsec:cech-descent}

\begin{theorem}[\v{C}ech descent for dg categories;
\ClaimStatusProvedElsewhere]
\label{thm:cech-descent}
\index{Cech descent!dg categories}
Cover $K3 = U_0 \cup U_1$, where $U_0 = K3 \setminus D$
(complement of a smooth divisor $D \in |H|$) and $U_1$ is a
tubular neighbourhood of $D$. The derived category $D^b(K3)$
is recovered as the homotopy fibre product
\[
D^b(K3) \;\simeq\; D^b(U_0) \times^h_{D^b(U_{01})} D^b(U_1),
\]
and the descent spectral sequence for Hochschild cohomology
\[
E_1^{p,q} = H^q\bigl(C^p(U_\bullet, \mathrm{HH}^*)\bigr)
\;\Longrightarrow\; \mathrm{HH}^{p+q}(D^b(K3))
\]
converges. The obstruction to gluing lives in
$H^2(K3, \cO_{K3}) \cong \bC$ (the Brauer class), reflecting
$h^{0,2}(K3) = 1$.
\end{theorem}

\subsection{Wall-crossing during chart transitions}
\label{subsec:wallcrossing-charts}

\begin{remark}[Cluster mutations and Bridgeland stability]
\label{rem:cluster-bridgeland}
\index{cluster mutation!chart transition}
\index{Bridgeland stability}
Mutation at a vertex $k$ of the McKay quiver produces a new
quiver $Q' = \mu_k(Q)$ with exchange matrix $B'$ given by the
Fomin--Zelevinsky rule. For $A_n$ quivers, the exchange graph
has $C_{n+1} = \binom{2n}{n}/(n+1)$ vertices (the $(n+1)$-th
Catalan number), corresponding to triangulations of an
$(n+3)$-gon. The Kontsevich--Soibelman wall-crossing formula
governs the change of DT invariants across stability walls.
For $K3$ with primitive Mukai vector $v$: the moduli space
$M_H(v) \cong \mathrm{Hilb}^n(K3)$ (Beauville) and the
DT transformation satisfies $\tau^{n+3} = [2]$
(Seidel--Thomas).
\end{remark}


%% ===============================================================
%% THE K3 CHIRAL ALGEBRA
%% ===============================================================

\section{The K3 chiral algebra}
\label{sec:k3-chiral-algebra}
\index{K3 surface!chiral algebra|textbf}
\index{N=4 superconformal algebra|textbf}

The K3 surface is the fundamental test case for the CY-to-chiral correspondence programme at $d = 2$.  Every CY$_2$ category is either a K3 category, an abelian surface category, or a quotient thereof, so the K3 chiral algebra is the prototype from which the entire $d = 2$ theory is built.  The K3 sigma model produces a $c = 6$ vertex algebra with $\cN = 4$ superconformal symmetry (the maximal supersymmetry at $c = 6$), and the enhanced symmetry constrains $\kappa_{\mathrm{ch}}$ below the naive Virasoro prediction.  The complete algebraic data (OPE structure, bar complex, shadow tower, Koszul dual, Mathieu moonshine) must be assembled and verified against two independent predictions: the $d=2$ member of the CY-to-chiral correspondence programme and the $\kappa_\bullet$-spectrum of the kappa-polysemy.

\subsection{The $\cN = 4$ superconformal algebra at $c = 6$}
\label{subsec:n4-sca}

\begin{definition}[Small $\cN = 4$ SCA at $c = 6$]
\label{def:n4-sca-c6}
\index{N=4 superconformal algebra!c=6}
The \emph{small $\cN = 4$ superconformal algebra} at central
charge $c = 6$ (equivalently, $\mathrm{SU}(2)_R$ level
$k_R = 1$, since $c = 6k_R$) has eight primary generators:
\begin{center}
\small
\renewcommand{\arraystretch}{1.2}
\begin{tabular}{llcl}
\textbf{Generator} & \textbf{Weight $h$} & \textbf{Parity}
 & $J^3$ \textbf{charge} \\
\hline
$T$ & $2$ & bosonic & $0$ \\
$G^+, G^-$ & $3/2$ & fermionic & $\pm 1/2$ \\
$\tilde{G}^+, \tilde{G}^-$ & $3/2$ & fermionic & $\mp 1/2$ \\
$J^{++}, J^{--}$ & $1$ & bosonic & $\pm 1$ \\
$J^3$ & $1$ & bosonic & $0$
\end{tabular}
\end{center}
The OPE structure at $c = 6$, $k_R = 1$:
\begin{align}
T(z)\,T(w) &\sim \frac{3}{(z-w)^4} + \frac{2T}{(z-w)^2}
 + \frac{\partial T}{z-w}, \label{eq:n4-TT} \\
J^3(z)\,J^3(w) &\sim \frac{1/2}{(z-w)^2}, \label{eq:n4-J3J3} \\
J^{++}(z)\,J^{--}(w) &\sim \frac{1}{(z-w)^2}
 + \frac{2J^3}{z-w}, \label{eq:n4-JppJmm} \\
G^+(z)\,G^-(w) &\sim \frac{2}{(z-w)^3} + \frac{2J^3}{(z-w)^2}
 + \frac{T + \partial J^3}{z-w}. \label{eq:n4-GpGm}
\end{align}
The remaining independent OPE relations at $c = 6$, $k_R = 1$
\textup{(}from the $64$-pair table of \texttt{cy\_n4sca\_k3\_engine.py}\textup{)}:
\begin{align}
J^3(z)\,G^\pm(w) &\sim \frac{\pm\tfrac{1}{2}\,G^\pm}{z-w}, \label{eq:n4-J3Gpm} \\
J^{++}(z)\,G^-(w) &\sim \frac{G^+}{z-w},\qquad
J^{--}(z)\,G^+(w) \sim \frac{G^-}{z-w}, \label{eq:n4-JpmGmp} \\
T(z)\,G^\pm(w) &\sim \frac{\tfrac{3}{2}\,G^\pm}{(z-w)^2}
 + \frac{\partial G^\pm}{z-w}, \label{eq:n4-TG} \\
T(z)\,J^a(w) &\sim \frac{J^a}{(z-w)^2} + \frac{\partial J^a}{z-w},
 \label{eq:n4-TJ} \\
\tilde{G}^+(z)\,\tilde{G}^-(w) &\sim \frac{2}{(z-w)^3}
 - \frac{2J^3}{(z-w)^2}
 + \frac{T - \partial J^3}{z-w}, \label{eq:n4-GtpGtm} \\
G^+(z)\,\tilde{G}^+(w) &\sim \frac{2\,J^{++}}{z-w},\qquad
G^-(z)\,\tilde{G}^-(w) \sim \frac{-2\,J^{--}}{z-w}.
 \label{eq:n4-GGt-cross}
\end{align}
All boson--boson and fermion--fermion OPEs not listed above
are either trivially determined by $\mathrm{SU}(2)_R$ symmetry
or regular \textup{(}no singular terms\textup{)}.
\end{definition}

\begin{proposition}[Modular characteristic of the K3 sigma model]
\label{prop:kappa-k3}
\ClaimStatusProvedHere
\index{modular characteristic!K3}
The modular characteristic of the K3 sigma model VOA is
\begin{equation}\label{eq:kappa-k3}
\kappa_{\mathrm{ch}}(\cA_{K3}) = 2 = \dim_\bC(K3).
\end{equation}
This is verified by six independent paths:
\begin{enumerate}[label=\textup{(\roman*)}]
\item \emph{Geometric.} For a CY $d$-fold sigma model:
 $\kappa_{\mathrm{ch}} = d$ \textup{(}index-theoretic computation via
 the Chern character of the chiral de~Rham complex\textup{)}.
\item \emph{Character.} $F_1 = \kappa_{\mathrm{ch}}/24 = 1/12$ matches the
 genus-$1$ anomaly of the $c = 6$ sigma model.
\item \emph{Complementarity.} $\kappa_{\mathrm{ch}}(\cA) + \kappa_{\mathrm{ch}}(\cA^!)
 = 0$ gives $\kappa_{\mathrm{ch}}^! = -2$.
\item \emph{Kummer orbifold.} At the Kummer point $K3 = T^4/\bZ_2$:
 the orbifold preserves $\kappa_{\mathrm{ch}} = 2$ from the smooth sigma model.
\item \emph{Gepner consistency.} The Gepner model $(2)^4$ gives
 $\kappa_{\mathrm{Gepner}} = 4 \times (3/4) = 3$ using the
 Virasoro formula $\kappa_{\mathrm{ch}} = c/2$ for each $\cN = 2$ factor.
 This is a DIFFERENT algebra from the $\cN = 4$ SCA
 \textup{(}$\kappa_{\mathrm{ch}}$ depends on the full algebra,
 not just the Virasoro subalgebra\textup{)}.
\item \emph{Additivity.} $\kappa_{\mathrm{ch}}(K3 \times E) =
 \kappa_{\mathrm{ch}}(K3) + \kappa_{\mathrm{ch}}(E) = 2 + 1 = 3 = \dim_\bC(K3 \times E)$
 \textup{(}Corollary~\textup{\ref{cor:shadow-extraction})}.
\end{enumerate}
\end{proposition}

\begin{proof}
Paths (i)--(iii) follow from Vol~I, Theorem~D and the index-theoretic
computation $\kappa_{\mathrm{ch}}(\Omega^{\mathrm{ch}}(\mathrm{CY}_d)) = d$
(the Chern character of the chiral de~Rham complex on a Ricci-flat
manifold reduces the genus-$1$ obstruction to $\chi(M)/12 = 2$
for K3 via $\chi(K3) = 24$). Path (iv): the Kummer orbifold
$T^4/\bZ_2$ has the same modular characteristic as the smooth K3
because $\kappa_{\mathrm{ch}}$ is a deformation invariant (it depends on the
chiral algebra structure, which is constant across K3 moduli by
curve independence). Path (v): the Gepner model uses the Virasoro
stress tensor for each $c = 3/2$ factor, giving $\kappa_{\mathrm{ch}} = c/2 = 3/4$
per factor; the physical sigma model has the full $\cN = 4$
structure whose $\kappa_{\mathrm{ch}} = 2$ is lower due to supersymmetric
Ward identities. Path (vi) follows from the additivity of $\kappa_{\mathrm{ch}}$
under tensor products
(Vol~I, Proposition~\ref{prop:independent-sum-factorization}).
\end{proof}

\begin{remark}[$\kappa_{\mathrm{ch}}(K3) = 2 \neq c/2 = 3$: modular characteristic vs central charge]
\label{rem:ap48-k3}
\index{modular characteristic!K3 example}
The K3 sigma model provides a sharp illustration: the
modular characteristic $\kappa_{\mathrm{ch}}(\cA_{K3}) = 2$ differs from the
Virasoro formula $c/2 = 3$. The reduction arises because the
$\cN = 4$ superconformal Ward identities constrain the genus-$1$
obstruction below what the Virasoro subalgebra alone would produce.
The naive Virasoro computation $\kappa_{\mathrm{ch}}(\mathrm{Vir}_6) = 3$ counts
the Virasoro contribution, but the full $\cN = 4$ algebra at $c = 6$
has $\kappa_{\mathrm{ch}} = 2k_R = 2$.
\end{remark}

\subsection{Bar complex of the $\cN = 4$ SCA}
\label{subsec:bar-n4}

\begin{proposition}[Bar complex structure; \ClaimStatusProvedHere]
\label{prop:bar-n4-sca}
\index{bar complex!N=4 SCA}
The bar complex $B(\cA_{K3})$ of the small $\cN = 4$ SCA at
$c = 6$ has:
\begin{enumerate}[label=\textup{(\roman*)}]
\item $B^1(\cA_{K3}) = \mathrm{span}\{s^{-1}a : a \text{ a
 generator}\}$ with $\dim B^1 = 8$
 \textup{(}one desuspended generator per primary;
 $|s^{-1}v| = |v| - 1$\textup{)}.
\item The bar differential $d\colon B^2 \to B^1$ is a
 rank-$16$ linear map extracting all singular OPE modes via
 the logarithmic kernel $d\log(z-w)$ \textup{(}the $d\log$
 measure absorbs one power of $(z-w)$, so bar pole orders are
 one less than OPE pole orders; the bar differential extracts
 all singular modes, not just the simple pole\textup{)}.
\item The algebra is chirally Koszul
 \textup{(}Vol~I, Proposition~\textup{\ref{prop:pbw-universality}}):
 the $\cN = 4$ SCA is freely strongly generated \textup{(}no null
 vectors in the universal vacuum module\textup{)}, so bar
 cohomology $H^*(B(\cA))$ is concentrated in bar degree~$1$.
\end{enumerate}
\end{proposition}

\subsection{Shadow depth classification}
\label{subsec:shadow-depth-k3}

\begin{proposition}[K3 sigma model: class M; \ClaimStatusProvedHere]
\label{prop:shadow-class-k3}
\index{shadow depth!K3}
The K3 sigma model VOA has shadow depth class~M
\textup{(}infinite tower, $r_{\max} = \infty$\textup{)}.
The shadow metric
$Q_L(t) = (2\kappa_{\mathrm{ch}} + 3\alpha t)^2 + 2\Delta\,t^2$ has
critical discriminant
\[
\Delta \;=\; 8\kappa_{\mathrm{ch}}\,S_4 \;\neq\; 0,
\]
so $Q_L$ is irreducible and the shadow tower does not terminate.
The Virasoro sector at $c = 6$ contributes
$Q^{\mathrm{contact}} = 10/[c(5c+22)] = 5/156 \neq 0$, and the
$\mathrm{SU}(2)_R$ sector at level $k_R = 1$ contributes nonzero
cubic shadow $S_3$ from the triple current OPE.
\end{proposition}

\begin{computation}[Shadow tower of the K3 sigma model through degree~$12$;
\ClaimStatusProvedHere]
\label{comp:shadow-tower-k3}
\index{shadow tower!K3 sigma model}
The MC recursion (Vol~I, Theorem~\textup{\ref{thm:riccati-algebraicity}})
with initial data $(\kappa_{\mathrm{ch}}, S_3, S_4) = (2, 2, 5/156)$ produces the
full shadow tower. Every value is an exact rational, verified to
satisfy the MC equation
$D_{\cA}(\Theta) + \tfrac{1}{2}[\Theta, \Theta] = 0$
at each degree.
\begin{center}
\renewcommand{\arraystretch}{1.2}
\begin{tabular}{rcl}
 \toprule
 $r$ & $S_r(\cA_\sigma)$ & Sign \\
 \midrule
 2 & $2$ & $+$ \\
 3 & $2$ & $+$ \\
 4 & $5/156$ & $+$ \\
 5 & $-1/26$ & $-$ \\
 6 & $3485/73008$ & $+$ \\
 7 & $-1145/18928$ & $-$ \\
 8 & $1179485/15185664$ & $+$ \\
 9 & $-1144885/11389248$ & $-$ \\
 10 & $309513983/2368963584$ & $+$ \\
 11 & $-1477140565/8686199808$ & $-$ \\
 12 & $490338857615/2217349914624$ & $+$ \\
 \bottomrule
\end{tabular}
\end{center}
Sign alternation begins at degree~$5$; magnitudes are bounded and
slowly decaying. The critical discriminant
$\Delta = 8\kappa_{\mathrm{ch}} S_4 = 20/39 \neq 0$ forces all $S_r \neq 0$ for
$r \geq 5$ by the single-line dichotomy theorem
(Vol~I, Theorem~\textup{\ref{thm:riccati-algebraicity}}),
confirming class~M with $r_{\max} = \infty$.

\smallskip\noindent\textit{Verification}:
\texttt{cy\_n4sca\_k3\_engine.py},
\texttt{cy\_bar\_n4sca\_engine.py}.
\end{computation}

\subsection{The K3 elliptic genus and Mathieu moonshine}
\label{subsec:k3-elliptic-genus}

\begin{definition}[K3 elliptic genus]
\label{def:k3-elliptic-genus}
\index{elliptic genus!K3|textbf}
\index{Jacobi form!phi_{0,1}}
The elliptic genus of $K3$ is
\begin{equation}\label{eq:k3-elliptic-genus}
Z_{K3}(\tau,z) \;=\; 2\,\phi_{0,1}(\tau,z),
\end{equation}
where $\phi_{0,1}$ is the unique weak Jacobi form of weight~$0$,
index~$1$ in the Eichler--Zagier normalization
$\phi_{0,1}(\tau,0) = 12$. At $z = 0$:
$Z_{K3}(\tau,0) = 24 = \chi(K3)$.
The Fourier expansion at $q^0$ gives
$Z_{K3}|_{q^0} = 2y + 20 + 2y^{-1}$,
the $\chi_y$-genus of K3 with $h^{1,1} = 20$,
$h^{2,0} = h^{0,2} = 1$.
\end{definition}

\begin{theorem}[Mathieu moonshine; \ClaimStatusProvedElsewhere]
\label{thm:mathieu-moonshine}
\index{Mathieu moonshine|textbf}
\index{M24|textbf}
The $\cN = 4$ decomposition of the K3 elliptic genus,
\begin{equation}\label{eq:mathieu-decomposition}
Z_{K3} \;=\; 20 \cdot \mathrm{ch}_{\mathrm{massless}}
\;+\; \sum_{n \geq 1} A_n \cdot \mathrm{ch}_{\mathrm{massive},n},
\end{equation}
where $\mathrm{ch}_{\mathrm{massless}}$ and
$\mathrm{ch}_{\mathrm{massive},n}$ are $\cN = 4$ characters at
$c = 6$, has coefficients $A_n$ that are dimensions of
representations of the Mathieu group $M_{24}$
\textup{(}Eguchi--Ooguri--Tachikawa 2010, proved by Gannon
2016\textup{)}:
\begin{center}
\small
\renewcommand{\arraystretch}{1.2}
\begin{tabular}{ccccccc}
$n$ & $1$ & $2$ & $3$ & $4$ & $5$ & $6$ \\
\hline
$A_n$ & $90$ & $462$ & $1540$ & $4554$ & $11592$ & $27830$
\end{tabular}
\end{center}
The $20$ massless states decompose as
$23_{\mathrm{std}} - 3_{\mathrm{triv}}$ under $M_{24}$.
\end{theorem}

\begin{remark}[Mock modular form from K3]
\label{rem:mock-modular-k3}
\index{mock modular form!K3}
The massive multiplicities assemble into the mock modular form
\begin{equation}\label{eq:mock-modular-h}
h(\tau) \;=\; 2q^{-1/8}\bigl(-1 + 45q + 231q^2
 + 770q^3 + 2277q^4 + \cdots\bigr),
\end{equation}
where $A_n^{\mathrm{full}} = 2 \cdot a_n$ with $a_n = 45, 231,
770,\ldots$ being the half-multiplicities. The shadow of $h$ in
the sense of Zwegers is $S(\tau) = 24\,\eta(\tau)^3$, a
weight-$3/2$ cusp form. The constant term $a_0 = -1$ gives
$A_0^{\mathrm{full}} = -2 = -\kappa_{\mathrm{ch}}(\cA_{K3})$: the modular
characteristic appears as the leading mock modular coefficient.

The ``shadow'' in mock modularity (Zwegers) and the ``shadow''
in the shadow obstruction tower (bar complex) are a priori
different mathematical objects. Their structural parallel
(both encode the failure of modular invariance) is a
connection to be explored, not an identification to be assumed.

The mock modular form $h(\tau)$ admits a canonical decomposition
via the Appell--Lerch sum. Let
$\mu(\tau,z) = \frac{-iy^{1/2}}{\vartheta_1(\tau,z)}
\sum_{n \in \bZ} \frac{(-1)^n q^{n(n+1)/2} y^n}{1 - q^n y}$
denote the Zwegers $\mu$-function. Then
\begin{equation}\label{eq:mock-modular-decomposition}
Z_{K3}(\tau,z)
\;=\;
\bigl(h(\tau) - 24\,\mu(\tau,z)\bigr)
\cdot \frac{\vartheta_1(\tau,z)^2}{\eta(\tau)^3},
\end{equation}
where the $24 = \chi(K3)$ multiplying $\mu$ is the Witten
index and the theta-to-eta factor
$\vartheta_1^2/\eta^3$ carries weight~$1/2$ and index~$1$.
The decomposition is verified to machine
precision (relative error $\sim 5.7 \times 10^{-16}$).
The polar term $h|_{q^{-1/8}} = -2 = -\kappa_{\mathrm{ch}}(\cA_{K3})$
reappears on the right-hand side as a consequence of the
Appell--Lerch cancellation.
\end{remark}

\begin{remark}[The factor $2$ in $Z_{K3} = 2\,\phi_{0,1}$ is $\kappa_{\mathrm{ch}}$]
\label{rem:factor-2-is-kappa}
\index{modular characteristic!as elliptic genus coefficient}
\index{Mathieu moonshine!bar complex coefficient}
The elliptic genus $Z_{K3} = 2\,\phi_{0,1}$ has an overall
factor of~$2$. This factor IS the modular characteristic:
$\kappa_{\mathrm{ch}}(\cA_{K3}) = 2$. The bar complex provides the
universal coefficient
\begin{equation}\label{eq:moonshine-multiplier}
A_n \;=\; \kappa_{\mathrm{ch}}(\cA_{K3}) \cdot \dim(\rho_n),
\end{equation}
where $\rho_n$ is the $M_{24}$-representation at level~$n$:
$\dim(\rho_1) = 45$, $\dim(\rho_2) = 231$,
$\dim(\rho_3) = 770$, $\dim(\rho_4) = 2277$. The
half-multiplicities $a_n = 45, 231, 770, \ldots$ are the
$M_{24}$ irreducible dimensions; the full multiplicities
$A_n = 2\,a_n = 90, 462, 1540, \ldots$ acquire the factor
$\kappa_{\mathrm{ch}} = 2$ from the bar complex.

The bar complex does NOT detect $M_{24}$ directly: it
detects $\kappa_{\mathrm{ch}} = 2$ at degree~$2$ and the cubic/quartic
shadow invariants at higher degrees. The $M_{24}$
representation structure is an \emph{internal} symmetry of
the massive BPS sector, invisible to the single-copy bar
construction. The bar complex provides the universal
multiplier; the group theory is orthogonal. The polar
coefficient $-2 = -\kappa_{\mathrm{ch}}$ in the mock modular
form~\eqref{eq:mock-modular-h} is the same $\kappa_{\mathrm{ch}}$
from a different direction.
\end{remark}

\begin{remark}[Twining genera and $M_{24}$ conjugacy classes]
\label{rem:twining-genera}
\index{twining genera}
For each conjugacy class $[g]$ of $M_{24}$, the twining genus
$Z_g(\tau,z) = \mathrm{Tr}_{\mathrm{RR}}\bigl(g \cdot
q^{L_0 - c/24}\,\bar{y}^{J_0}\,(-1)^{F+\bar{F}}\bigr)$ is a
weak Jacobi form of weight~$0$, index~$1$ for
$\Gamma_0(N_g)$, where $N_g = \mathrm{ord}(g)$. Of the $26$
conjugacy classes of $M_{24}$, only $21$ appear as symmetries of
K3 sigma models (Gaberdiel--Hohenegger--Volpato 2010/2011): the
five missing classes ($7A, 7B, 15A, 15B, 23A/B$) have Frame
shapes incompatible with the K3 lattice. Each realized class $g$
has a Frame shape $\pi_g = \prod i^{a_i}$ encoding the permutation
action on $24$ letters, with associated eta product
$\eta_g(\tau) = \prod_i \eta(i\tau)^{a_i}$.
The $21$ realized classes with their Frame shapes and orders:
\begin{center}
\scriptsize
\renewcommand{\arraystretch}{1.15}
\begin{tabular}{llr|llr|llr}
Class & Frame shape & $N_g$
& Class & Frame shape & $N_g$
& Class & Frame shape & $N_g$ \\
\hline
$1A$ & $1^{24}$ & $1$
& $2A$ & $1^82^8$ & $2$
& $2B$ & $2^{12}$ & $2$ \\
$3A$ & $1^63^6$ & $3$
& $4A$ & $1^42^24^4$ & $4$
& $4B$ & $2^44^4$ & $4$ \\
$4C$ & $4^6$ & $4$
& $5A$ & $1^45^4$ & $5$
& $6A$ & $1^22^13^26^2$ & $6$ \\
$6B$ & $1^22^23^16^1$ & $6$
& $8A$ & $1^22^14^18^2$ & $8$
& $10A$ & $1^22^15^110^1$ & $10$ \\
$11A$ & $1^211^2$ & $11$
& $12A$ & $1^12^14^13^16^112^1$ & $12$
& $12B$ & $2^14^16^112^1$ & $12$ \\
$14A$ & $1^12^17^114^1$ & $14$
& $14B$ & $1^12^17^114^1$ & $14$
& $21A$ & $1^13^17^121^1$ & $21$ \\
$21B$ & $1^13^17^121^1$ & $21$
& $23A$ & $1^123^1$ & $23$
& $23B$ & $1^123^1$ & $23$
\end{tabular}
\end{center}
The $23A$ and $23B$ classes are listed for completeness but are
among the five NOT realized as K3 symmetries. The Frame shape
determines the eta product $\eta_g$ and hence the shadow of the
twined mock modular form; the shadow of $h_g$ is proportional
to $\eta_g$.
\end{remark}

\subsection{Mathieu moonshine through the K3 Yangian}
\label{subsec:mathieu-yangian}
\index{Mathieu moonshine!K3 Yangian}
\index{Yangian!K3!Mathieu moonshine}

The K3 Yangian $Y(\frakg_{K3})$ of
Section~\ref{sec:ade-yangian-level1} provides a structural
framework for understanding WHY $M_{24}$ acts on the K3
elliptic genus.

\begin{conjecture}[$M_{24}$ action on $Y(\frakg_{K3})$ representations]
\label{conj:m24-yangian-action}
\ClaimStatusConjectured
\index{M24@$M_{24}$!Yangian action}
The Mathieu group $M_{24}$ acts on the representation category
$\mathrm{Rep}(Y(\frakg_{K3}))$ via its natural permutation
of the $24$~Mukai lattice directions.  Concretely:
\begin{enumerate}[label=\textup{(\roman*)}]
\item
The $24$ Yangian parameters $h_1, \ldots, h_{24}$ live
in the complexified Mukai lattice
$\widetilde{H}(K3, \bC) \cong \bC^{24}$,
subject to the CY$_2$ constraint $\sum h_i = 0$.
The group $M_{24}$ permutes the $h_i$, preserving the
constraint (since permutations preserve sums).
\item
The structure function
$g_{K3}(z) = \prod_{i=1}^{24}(z - h_i)/(z + h_i)$
is a symmetric function of the $h_i$, hence invariant
under $M_{24}$.  The spectral parameter~$z$ is a Yangian
shift parameter {\upshape(}$z$ spectral, not worldsheet;
Remark~\ref{rem:z-spectral-vs-worldsheet}{\upshape)},
invisible to the $M_{24}$ action.
\item
The charge-$1$ representation $V = \bC^{24}$ decomposes
under $M_{24}$ as $23_{\mathrm{std}} \oplus 1_{\mathrm{triv}}$,
matching the constraint hyperplane $\sum h_i = 0$ (dimension~$23$)
plus the normal direction (dimension~$1$).
\end{enumerate}
The conjecture is conditional on the existence of
$Y(\frakg_{K3})$.
\end{conjecture}

\begin{remark}[Twined structure functions and Frame shapes]
\label{rem:twined-structure-fn-yangian}
\index{Frame shape!twined structure function}
For each conjugacy class $[g]$ of $M_{24}$ with Frame
shape $\pi_g = \prod i^{a_i}$, the twined
R-matrix trace restricts the structure function product
to the $a_1$~fixed-point directions, yielding a twined
structure function of degree $(a_1, a_1)$:
\[
g_g^{\mathrm{twined}}(z)
\;=\;
\prod_{i:\, g(i) = i} \frac{z - h_i}{z + h_i},
\qquad
\deg = (a_1, a_1).
\]
For $1A$: degree $(24,24)$.  For $2A$ ($a_1 = 8$):
degree $(8,8)$.  For $2B$ ($a_1 = 0$): trivially~$1$.

The twined bar Euler product decomposes as the eta product
$\eta_g(\tau) = \prod_{i:\,a_i \neq 0} \eta(i\tau)^{a_i}$,
consistent with the shadow of the twined mock modular
form~$h_g$ (Remark~\ref{rem:twining-genera}).
\end{remark}

\begin{remark}[Orthogonality of bar complex and group theory]
\label{rem:bar-m24-orthogonality}
\index{Mathieu moonshine!orthogonality principle}
The moonshine decomposition
$A_n = \kappa_{\mathrm{ch}} \cdot \dim(\rho_n)$
(equation~\eqref{eq:moonshine-multiplier})
exhibits a factorization into two independent contributions:
the bar complex provides $\kappa_{\mathrm{ch}} = 2$ (a
topological invariant of the K3 chiral algebra, computed from the
shadow tower); the group theory provides $\dim(\rho_n)$ (an
internal symmetry of the massive BPS sector, invisible to the
single-copy bar construction).  Neither factor determines the
other.  The K3 Yangian is the proposed structural home for
this factorization: $\kappa_{\mathrm{ch}}$ comes from the bar
complex of $Y(\frakg_{K3})$, while $\dim(\rho_n)$ comes from
the $M_{24}$ action on $\mathrm{Rep}(Y(\frakg_{K3}))$.
\end{remark}

\begin{remark}[Umbral moonshine and the Niemeier Yangian landscape]
\label{rem:umbral-yangian}
\index{Umbral moonshine!Yangian}
\index{Niemeier lattice!Yangian specialization}
The $23$ Niemeier lattices with nonempty root system~$R(N)$
each determine an Umbral moonshine module
(Cheng--Duncan--Harvey, arXiv:1204.2779).
The Umbral group $G_N$ and lambency~$\ell$ (the Coxeter number
of~$R(N)$) parametrize the mock modular forms whose
coefficients are $G_N$-representation dimensions.  For the
$A_1^{24}$ Niemeier lattice: $G_N = M_{24}$, $\ell = 2$,
recovering the original Eguchi--Ooguri--Tachikawa observation.

At each Niemeier maximal enhancement point of the K3 moduli
space (Proposition~\ref{prop:niemeier-shadow-landscape}),
the K3 Yangian parameters $h_i$ specialize to incorporate the
root structure of~$N$.  The Umbral group~$G_N$ acts on the
Yangian representation category at this specialization.
The shadow data is identical for all~$24$ Niemeier lattices
($\kappa_{\mathrm{ch}} = 24$, class~$G$, depth~$2$ for the
lattice VOA), so the Umbral group is invisible to the shadow
tower; it is detected only through the representation theory
of the specialized Yangian.
\end{remark}

\noindent\textit{Verification}: $50$ tests in
\texttt{mathieu\_moonshine\_yangian.py}:
$25$~conjugacy classes (Frame shapes summing to~$24$, cycle
lengths dividing order, trace consistency); $9$~moonshine
coefficients ($A_n = 2\,a_n$); twined eta products for
$1A$, $2A$, $2B$, $3A$ matching Ramanujan tau and eta product
expansions; $M_{24}$ action on charge-$1$ and
$\mathrm{Sym}^2$ representations;
$23$~Niemeier root systems with Umbral data;
trace-of-power formula for all classes and powers.

\subsection{Lattice VOA and Gepner model}
\label{subsec:lattice-gepner}

\begin{remark}[Lattice VOA from $H^2(K3,\bZ)$]
\label{rem:lattice-voa-k3}
\index{lattice VOA!K3}
\index{K3 lattice}
The K3 cohomology lattice $L_{K3} = H^2(K3,\bZ) \cong U^3 \oplus
(-E_8)^2$ has rank~$22$, signature $(3,19)$, and discriminant
$\det(G) = -1$. Since $L_{K3}$ is indefinite, the naive theta
function diverges. The negative-definite sublattice $(-E_8)^2$
of rank~$16$ produces a well-defined lattice VOA
$V_{(-E_8)^2}$ with $c = 16$ and $\kappa_{\mathrm{ch}} = 16$
($\kappa_{\mathrm{ch}} = \mathrm{rank}$ for lattice VOAs). The
theta series $\Theta_{E_8}(\tau) = E_4(\tau) = 1 + 240q + \cdots$
(the weight-$4$ Eisenstein series, since $E_8$ is the unique even
unimodular lattice of rank~$8$).

The full Mukai lattice $\widetilde{H}(K3,\bZ) = U^4 \oplus
(-E_8)^2$ has rank~$24$, signature $(4,20)$, and is the lattice
relevant for Bridgeland stability conditions and derived
categories.
\end{remark}

\begin{remark}[Gepner model $(2)^4$]
\label{rem:gepner-model-k3}
\index{Gepner model|textbf}
The Gepner model for the quartic K3 (Fermat quartic
$x_1^4 + x_2^4 + x_3^4 + x_4^4 = 0$ in $\bC\bP^3$) is the
tensor product of four copies of the $\cN = 2$ minimal model at
level $k = 2$ ($c = 3k/(k+2) = 3/2$), with GSO
projection and $\bZ_4$ orbifold. Total $c = 4 \times 3/2 = 6$.
The chiral ring is isomorphic to the Jacobian ring
$\mathrm{Jac}(W) = \bC[x_1,\ldots,x_4]/(\partial W/\partial x_i)$
with $\dim \mathrm{Jac}(W) = 3^4 = 81$ before orbifold. The
symmetry group at the Gepner point contains
$(\bZ/4\bZ)^4 \rtimes S_4$ (order $6144$) and embeds into
the Conway group $\mathrm{Co}_1$.
\end{remark}

\subsection{Chiral de~Rham complex on K3}
\label{subsec:cdr-k3}

\begin{remark}[CDR on K3]
\label{rem:cdr-k3}
\index{chiral de Rham complex!K3}
The chiral de~Rham complex $\Omega^{\mathrm{ch}}(K3)$ of
Malikov--Schechtman--Vaintrob is a sheaf of vertex
superalgebras on K3 with central charge $c = 0$ (the local
contributions $c_{\beta\gamma} = +2$ and $c_{bc} = -2$ per
complex dimension cancel globally). The CDR cohomology
recovers the K3 elliptic genus:
$\mathrm{ch}\bigl(H^*\bigl(K3, \Omega^{\mathrm{ch}}\bigr)\bigr)
= \mathrm{Ell}(K3; q, y) = 2\,\phi_{0,1}(\tau, z)$
(Borisov--Libgober).

On a hyperk\"ahler manifold (such as K3), the CDR carries
$\cN = 4$ superconformal symmetry at $c = 0$: the
topological $\cN = 4$ algebra, distinct from the physical
$\cN = 4$ at $c = 6$. Both give
$\kappa_{\mathrm{ch}} = \dim_\bC(K3) = 2$.
\end{remark}

\subsection{Koszul dual of the K3 chiral algebra}
\label{subsec:koszul-dual-k3}

\begin{proposition}[Koszul dual of $\cA_{K3}$; \ClaimStatusProvedHere]
\label{prop:koszul-dual-k3}
\index{Koszul dual!K3 chiral algebra}
The Koszul dual of the $\cN = 4$ SCA at $c = 6$, $k_R = 1$ has:
\begin{itemize}
\item $\kappa_{\mathrm{ch}}(\cA_{K3}^!) = -2$ \textup{(}by complementarity:
 $\kappa_{\mathrm{ch}} + \kappa_{\mathrm{ch}}' = 0$ for free-field/CY sigma models\textup{)}.
\item $c(\cA_{K3}^!) = -6$
 \textup{(}the $\cN = 4$ SCA at negative central charge, with
 $k_R' = -1$, since $\kappa_{\mathrm{ch}} = 2k_R$ and $\kappa_{\mathrm{ch}}' = -2$ gives
 $k_R' = -1$, hence $c' = 6k_R' = -6$\textup{)}.
\item $F_g(\cA^!) = -2 \cdot \lambda_g^{\mathrm{FP}}$ at all
 genera \textup{(all-weight, with cross-channel correction $\delta F_g^{\mathrm{cross}}$ at $g \geq 2$; Vol~I)}.
\end{itemize}
The Verdier intertwining \textup{(}Vol~I, Theorem~A;
Convention~\textup{\ref{conv:bar-coalgebra-identity}):}
$D_{\mathrm{Ran}}(B(\cA_{K3})) \simeq B(\cA_{K3}^!)$.
The derived centre $Z^{\mathrm{der}}_{\mathrm{ch}}(\cA_{K3})$
\textup{(}the universal bulk\textup{)} is a separate object
from $\cA^!$.
\end{proposition}


%% ===============================================================
%% ABELIAN (2,0) PUSHFORWARD AND THE KAPPA-COLLAPSE
%% ===============================================================

\subsection{The abelian $(2,0)$ pushforward}
\label{subsec:pushforward-k3}
\index{pushforward chiral algebra!K3}

For $X = S \times E$ with $S$ a K3 surface and $E$ an elliptic
curve, the derived pushforward $R\pi_*(\Omega^2_{X,\mathrm{cl}})$
along $\pi \colon S \times E \to E$ produces a chiral algebra on~$E$.

\begin{construction}[The pushforward chiral algebra]
\label{constr:pushforward-chiral-algebra}
\index{abelian (2,0) theory!pushforward}
The \emph{abelian $(2,0)$ pushforward chiral algebra} on $E$ is
\[
 \cA_{\mathrm{free}} \;=\; \pi_*\bigl(
 \mathrm{Obs}^q(\mathcal{T}_{(2,0)}|_{S \times E})\bigr),
\]
the pushforward of the quantum observables of the holomorphically
twisted abelian $(2,0)$ theory along $\pi \colon S \times E \to E$.
By the Costello--Gwilliam theorem (proper pushforward preserves
factorization), this is a factorization algebra on $\mathrm{Ran}(E)$,
hence a chiral algebra on~$E$.

The fiber cohomology $R\pi_*(\Omega^2_X)$ decomposes via
K\"unneth into contributions from $H^*(S)$:
\begin{center}
\small
\renewcommand{\arraystretch}{1.2}
\begin{tabular}{llc}
 \toprule
 \textbf{Hodge group of $S$} & \textbf{Sheaf on $E$}
 & \textbf{Dimension} \\
 \midrule
 $H^0(S, \Omega^2_S) = \bC\langle\sigma\rangle$
 & $\cO_E$ & $1$ \\
 $H^2(S, \Omega^2_S) = \bC$ (Serre dual)
 & $\cO_E$ & $1$ \\
 $H^1(S, \Omega^1_S) = \bC^{20}$
 & $\Omega^1_E$ & $20$ \\
 $H^0(S, \cO_S) = \bC$ & $\Omega^2_E \simeq \cO_E$ & $1$ \\
 $H^2(S, \cO_S) = \bC$ & $\Omega^2_E \simeq \cO_E$ & $1$ \\
 \midrule
 \multicolumn{2}{l}{Total before closure constraint} & $24$ \\
 \bottomrule
\end{tabular}
\end{center}
The closure condition modifies this count: a section
$\sigma \otimes f(z)$ of the $H^{2,0}(S) \otimes \cO_E$
summand is closed only when $d_E(f) = 0$, i.e.\ $f$ is constant.
The naive total of~$24$ is therefore an upper bound for the rank
of the pushforward of closed $2$-forms. The $H^{1,1}(S)$ sector
(contributing $20$ weight-$1$ currents) survives the closure
condition, as its members are closed along $S$ and the
$\Omega^1_E$ factor prevents a $d_E$-constraint.
\end{construction}

\begin{proposition}[OPE structure of the pushforward;
\ClaimStatusProvedHere]
\label{prop:pushforward-ope}
\index{pushforward chiral algebra!OPE}
The chiral algebra $\cA_{\mathrm{free}}$ is of Heisenberg/lattice
type. Its OPE is determined by the propagator of the
$6$-dimensional theory pushed forward along~$S$:
\begin{enumerate}[label=\textup{(\roman*)}]
 \item The $20$ currents $J_a(z)$ from
 $H^{1,1}(S) \otimes \Omega^1_E$ have double-pole OPE:
 \[
 J_a(z)\, J_b(w) \;\sim\; \frac{\eta_{ab}}{(z-w)^2},
 \]
 where $\eta_{ab} = \int_S \omega_a \wedge \omega_b$ is the
 intersection form restricted to $H^{1,1}(S)$, of signature
 $(1,19)$.
 \item The remaining generators \textup{(}from $H^{2,0}$,
 $H^{0,2}$, $H^0(\cO_S)$, $H^2(\cO_S)$\textup{)} have at most
 double-pole OPE.
 \item No poles of order $\geq 3$ appear: the $6$-dimensional
 theory is free with at most quadratic action.
\end{enumerate}
In particular, $\cA_{\mathrm{free}}$ is a lattice vertex algebra
$V_\Lambda$ with $\Lambda$ determined by the fiber cohomology
surviving the closure constraint. By the Vol~I lattice curvature
theorem, $\kappa_{\mathrm{ch}}(V_\Lambda) = \mathrm{rank}(\Lambda)$ for any
lattice VA, shadow class~$\mathbf{G}$
\textup{(}Gaussian, $r_{\max} = 2$\textup{)}, and the shadow
obstruction tower terminates:
$\Theta_{\cA_{\mathrm{free}}} = \kappa_{\mathrm{ch}} \cdot \omega_g$ at all genera
\textup{(}UNIFORM-WEIGHT; class~G, scalar formula exact\textup{)}.
\end{proposition}

\begin{remark}[The $\kappa_{\mathrm{ch}}$-collapse: rank $\to$ $\dim_\bC$]
\label{rem:kappa-collapse}
\index{kappa-collapse@$\kappa_{\mathrm{ch}}$-collapse}
The passage from the abelian $(2,0)$ pushforward to the K3 sigma
model exhibits a non-perturbative collapse of the modular
characteristic:
\[
 \kappa_{\mathrm{ch}}(\cA_{\mathrm{free}}) \;=\; \mathrm{rank}(\Lambda)
 \;\in\; \{22, 24\},
 \qquad
 \kappa_{\mathrm{ch}}(\cA_\sigma) \;=\; 2 \;=\; \dim_\bC(K3).
\]
This is not a small correction. No deformation of
$\cA_{\mathrm{free}}$ within the class of lattice vertex algebras
can produce $\kappa_{\mathrm{ch}} = 2$: the Vol~I lattice curvature theorem gives
$\kappa_{\mathrm{ch}}(V_\Lambda) = \mathrm{rank}(\Lambda)$ for all lattice VAs
independent of deformation parameters.
The $\cN = 4$ supersymmetry constrains the genus-$1$ curvature to
equal the complex dimension of the target, regardless of the
free-field rank.

\medskip
The three levels of structure are:
\begin{center}
\renewcommand{\arraystretch}{1.3}
\begin{tabular}{llcll}
 \toprule
 \textbf{Object} & \textbf{Symbol} & $\kappa_\bullet$ & \textbf{Class}
 & \textbf{Depth} \\
 \midrule
 Abelian $(2,0)$ pushforward & $\cA_{\mathrm{free}}$
 & $\kappa_{\mathrm{ch}} = \mathrm{rank}(\Lambda) \in \{22, 24\}$
 & $\mathbf{G}$ & $2$ \\
 K3 sigma model VOA & $\cA_\sigma$
 & $\kappa_{\mathrm{ch}} = 2$ & $\mathbf{M}$ & $\infty$ \\
 BKM global & $G(K3 \times E)$
 & $\kappa_{\mathrm{BKM}} = 5$ & $\mathbf{M}$ & $\infty$ \\
 \bottomrule
\end{tabular}
\end{center}
The chain $\mathrm{rank}(\Lambda) \to 2 \to 5$ reflects: (i)~the
$\cN = 4$ SCA constraining $\kappa_{\mathrm{ch}}$ from $\mathrm{rank}$ to
$\dim_\bC$ (introducing $S_3, S_4 \neq 0$ and upgrading the
shadow class from~$\mathbf{G}$ to~$\mathbf{M}$); (ii)~the passage
from the single-copy chiral algebra
($\kappa_{\mathrm{ch}} = 3 = \kappa_{\mathrm{ch}}(K3) + \kappa_{\mathrm{ch}}(E) = 2 + 1$,
by additivity) to the second-quantized BKM modular characteristic
($\kappa_{\mathrm{BKM}} = 5$).
\end{remark}


%% ===============================================================
%% K3 x E AS A CALABI--YAU THREEFOLD
%% ===============================================================

\section{The Calabi--Yau threefold \texorpdfstring{$K3 \times E$}{K3 x E}:
geometry and invariants}
\label{sec:k3xe-geometry}
\index{K3 x E!geometry|textbf}

\subsection{Hodge diamond and topological invariants}
\label{subsec:hodge-k3xe}

\begin{proposition}[Hodge diamond of $K3 \times E$;
\ClaimStatusProvedElsewhere]
\label{prop:hodge-k3xe}
\index{Hodge diamond!K3 x E}
The K\"unneth formula gives the Hodge diamond of
$K3 \times E$ \textup{(}$\dim_\bC = 3$\textup{)}:
\[
\renewcommand{\arraystretch}{1.1}
\begin{array}{ccccccc}
& & & 1 & & & \\
& & 1 & & 1 & & \\
& 1 & & 21 & & 1 & \\
1 & & 21 & & 21 & & 1 \\
& 1 & & 21 & & 1 & \\
& & 1 & & 1 & & \\
& & & 1 & & &
\end{array}
\]
Key invariants:
\begin{itemize}
\item $h^{1,1} = h^{2,1} = 21$, giving
 $\chi(K3 \times E) = 2(h^{1,1} - h^{2,1}) = 0$.
\item Betti numbers: $b_0 = b_6 = 1$, $b_1 = b_5 = 2$,
 $b_2 = b_4 = 23$, $b_3 = 44$. Total: $\sum b_k = 96$.
\item $K_0(K3 \times E) \cong \bZ^{48}$
 \textup{(}$K_0(K3) \otimes K_0(E) = \bZ^{24} \otimes \bZ^2$\textup{)}.
\end{itemize}
\end{proposition}

\begin{proof}
Apply K\"unneth: $h^{p,q}(K3 \times E) = \sum_{a+c=p,\; b+d=q}
h^{a,b}(K3) \cdot h^{c,d}(E)$. For example:
$h^{1,1} = h^{1,1}(K3) \cdot h^{0,0}(E) + h^{0,0}(K3) \cdot
h^{1,1}(E) = 20 + 1 = 21$. The alternating sum
$1 - 2 + 23 - 44 + 23 - 2 + 1 = 0$ verifies $\chi = 0$.
\end{proof}

\subsection{HKR decomposition and deformation space}
\label{subsec:hkr-k3xe}

\begin{proposition}[Deformation space of $D^b(K3 \times E)$;
\ClaimStatusProvedElsewhere]
\label{prop:hkr-k3xe}
\index{HKR decomposition!K3 x E}
\index{Hochschild cohomology!K3 x E}
By the HKR isomorphism for a smooth CY $3$-fold:
$\mathrm{HH}^n(X) = \bigoplus_{p+q=n} h^{3-p,\,q}(X)$.
For $K3 \times E$:
\begin{center}
\small
\renewcommand{\arraystretch}{1.2}
\begin{tabular}{lll}
$\mathrm{HH}^0 = 1$ & automorphisms & $h^{3,0}$ \\
$\mathrm{HH}^1 = 2$ & outer derivations & $h^{3,1} + h^{2,0}$ \\
$\mathrm{HH}^2 = 23$ & first-order deformations & $h^{3,2} + h^{2,1} + h^{1,0}$ \\
$\mathrm{HH}^3 = 44$ & obstructions & $h^{3,3} + h^{2,2} + h^{1,1} + h^{0,0}$
\end{tabular}
\end{center}
The $23$-dimensional deformation space: $21$ complex structure
\textup{(}from $h^{2,1}$\textup{)}, $1$ B-field \textup{(}from
$h^{1,0}$\textup{)}, $1$ volume \textup{(}from
$h^{3,2}$\textup{)}. Total
$\dim \mathrm{HH}^* = 96$.
\end{proposition}

\subsection{Noncommutative deformations and Brauer group}
\label{subsec:nc-brauer}

\begin{remark}[Brauer group of $K3 \times E$]
\label{rem:brauer-k3xe}
\index{Brauer group!K3 x E}
$\mathrm{Br}(K3 \times E) = \mathrm{Br}(K3) \oplus \mathrm{Br}(E)$
contains $(\bQ/\bZ)^{21}$ (K3 at Picard rank $\rho = 1$) plus
$\bQ/\bZ$ from $E$. The K3 symplectic form provides the unique
Poisson structure $\pi = \omega^{-1}$ (since $\dim H^0(\bigwedge^2
T_{K3}) = 1$), but the $E$-direction carries none
($\bigwedge^2 T_E = 0$).
\end{remark}

\subsection{Autoequivalences and semiorthogonal structure}
\label{subsec:sod-k3xe}

\begin{proposition}[Indecomposability of $D^b(K3 \times E)$;
\ClaimStatusProvedElsewhere]
\label{prop:sod-k3xe}
\index{semiorthogonal decomposition!K3 x E}
$D^b(K3)$ admits no proper semiorthogonal decomposition
\textup{(}Bridgeland, Kawamata, Huybrechts: $\omega_{K3} = \cO$
forces Serre functor $[2]$, and the only fractional CY category
with Serre $[2]$ is $D^b(K3)$ itself\textup{)}. Hence
$D^b(K3 \times E) = D^b(K3) \otimes D^b(E)$ is indecomposable.

Autoequivalences: $\mathrm{Aut}(D^b(K3)) \to
O^+(\widetilde{H}(K3,\bZ))$ via oriented Hodge isometries
(line bundle twists, shifts, spherical twists $T_E$ with
$T_E^2 = [-2]$). $\mathrm{Aut}(D^b(E)) =
(E \times E^\vee) \rtimes (\mathrm{SL}(2,\bZ) \rtimes \bZ)$.
The virtual dimension of all moduli problems on $K3 \times E$
vanishes: $\mathrm{vdim} = -\chi(\cE,\cE) = 0$.
\end{proposition}

\subsection{Factorization homology}
\label{subsec:fact-homology-k3xe}

\begin{remark}[Factorization homology on $K3 \times E$]
\label{rem:fact-homology-k3xe}
\index{factorization homology!K3 x E}
The factorization homology of the HT-twisted sigma model:
$\int_{K3 \times E} A^{\mathrm{HT}} \simeq
H^*(K3 \times E,\, \Omega^{\mathrm{ch}})$.
The CDR character gives $\chi(\Omega^{\mathrm{ch}}) = 0$, but the
Hodge-graded character is nontrivial.
$\dim \mathrm{HH}^*(D^b(K3 \times E)) = 96$.
\end{remark}


%% ===============================================================
%% ENUMERATIVE GEOMETRY
%% ===============================================================

\section{Enumerative geometry of \texorpdfstring{$K3 \times E$}{K3 x E}}
\label{sec:enum-geom-k3xe}
\index{enumerative geometry!K3 x E|textbf}

\subsection{Gromov--Witten theory and BPS invariants}
\label{subsec:gw-bps}

\begin{remark}[Reduced GW theory on K3]
\label{rem:gw-k3}
\index{Gromov--Witten!reduced|textbf}
\index{Yau--Zaslow formula}
For K3, the standard virtual class vanishes for $\beta \neq 0$
(cosection from $H^{2,0}(K3) = \bC$). The reduced theory
(Bryan--Leung, Maulik--Pandharipande--Thomas) removes this
obstruction. The KKV formula (proved by Pandharipande--Thomas
2016) gives the BPS state counts. The genus-$0$ BPS invariants
satisfy the Yau--Zaslow formula:
\begin{equation}\label{eq:yau-zaslow}
\sum_{h \geq 0} n^0_h\,q^h \;=\; \prod_{n \geq 1}
\frac{1}{(1 - q^n)^{24}}.
\end{equation}
\end{remark}

\begin{proposition}[BPS invariants of K3; \ClaimStatusProvedElsewhere]
\label{prop:bps-k3}
\index{BPS invariants!K3}
Selected BPS invariants \textup{(}all integers, proved by
Pandharipande--Thomas 2016\textup{)}:
\begin{center}
\small
\renewcommand{\arraystretch}{1.2}
\begin{tabular}{lcccccccc}
$h$ & $0$ & $1$ & $2$ & $3$ & $4$ & $5$ & $6$ & $7$ \\
\hline
$n^0_h$ & $1$ & $24$ & $324$ & $3200$ & $25650$ & $176256$
 & $1073720$ & $5930496$ \\
$n^1_h$ & $0$ & $-2$ & $-54$ & $-800$ & $-8550$ & $-73440$
 & $-536860$ & $-3464264$ \\
$n^2_h$ & $0$ & $0$ & $3$ & $104$ & $1701$ & $19656$
 & $176358$ & $1316808$ \\
$n^3_h$ & $0$ & $0$ & $0$ & $-4$ & $-155$ & $-2832$
 & $-34470$ & $-320320$
\end{tabular}
\end{center}
Signed convention: $n^g_h = (-1)^g |n^g_h|$. The genus-$0$
values are the Yau--Zaslow numbers \textup{(}OEIS A006922\textup{)};
$n^0_1 = 24 = \chi(K3)$.
The Hilbert scheme Euler characteristics
$\chi(\operatorname{Hilb}^n(K3))$ provide the
genus-$0$ column: $\chi(\operatorname{Hilb}^n) = n^0_n$,
with generating function
$\sum_{n \geq 0} \chi(\operatorname{Hilb}^n)\,q^n =
\prod_{k \geq 1}(1-q^k)^{-24}$
\textup{(}G\"ottsche formula\textup{)}. The first eight values
are $1, 24, 324, 3200, 25650, 176256, 1073720, 5930496$.
\end{proposition}

\subsection{Bridgeland stability on K3}
\label{subsec:bridgeland-k3}

\begin{remark}[Stability conditions]
\label{rem:bridgeland-k3}
\index{Bridgeland stability!K3}
$\mathrm{Stab}^\dagger(K3) \to P^+(K3)$ is a covering of the
period domain (Bridgeland 2008). Walls of marginal stability for
Mukai vectors $v = v_1 + v_2$ are semicircles in the
$(B,\omega)$-plane; crossing them applies the KS wall-crossing.
For primitive $v$ with $v^2 = 2n - 2$:
$\Omega(v) = \chi(\mathrm{Hilb}^n(K3))$.
\end{remark}


%% ===============================================================
%% BPS SPECTRUM AND BLACK HOLE PHYSICS
%% ===============================================================

\section{BPS spectrum and black hole entropy}
\label{sec:bps-bh-k3xe}
\index{BPS spectrum!K3 x E|textbf}
\index{black hole entropy!K3 x E}

\subsection{Charge lattice and BPS states}
\label{subsec:charge-lattice}

\begin{proposition}[Charge lattice; \ClaimStatusProvedElsewhere]
\label{prop:charge-lattice}
\index{charge lattice!K3 x E}
Type IIA on $K3 \times E$ gives $4$d $\cN = 4$ supergravity
with $22$ vector multiplets. The D-brane charge lattice
$\Gamma = H^{\mathrm{even}}(K3 \times E, \bZ)$:
\begin{center}
\small
\renewcommand{\arraystretch}{1.2}
\begin{tabular}{llr}
\textbf{Cohomology} & \textbf{Brane type} & \textbf{Rank} \\
\hline
$H^0 = \bZ$ & D6-branes & $1$ \\
$H^2 = \bZ^{23}$ & D4-branes & $23$ \\
$H^4 = \bZ^{23}$ & D2-branes & $23$ \\
$H^6 = \bZ$ & D0-branes & $1$
\end{tabular}
\end{center}
Total: $\mathrm{rank}\,\Gamma = 48$.
\end{proposition}

\subsection{Black hole entropy from $1/\Phi_{10}$}
\label{subsec:bh-entropy}

\begin{proposition}[BMPV and $\frac{1}{4}$-BPS entropy;
\ClaimStatusProvedElsewhere]
\label{prop:bmpv}
\index{BMPV black hole}
In $5$d \textup{(}IIA on $K3 \times S^1$\textup{)}: the BMPV black hole
with charges $(Q_1, Q_5, n)$ has
$S_{\mathrm{BH}} = 2\pi\sqrt{Q_1\,Q_5\,n}$
\textup{(}Cardy with $c_L = 6\,Q_1\,Q_5$\textup{)}.
In $4$d \textup{(}IIA on $K3 \times T^2$\textup{)}: the
$\frac{1}{4}$-BPS dyon with discriminant $\Delta$ has
\[
S_{\mathrm{BH}} = 4\pi\sqrt{\Delta},
\]
and the DVV formula $Z_{1/4\text{-BPS}} = 1/\Phi_{10}$
\textup{(}\eqref{eq:dvv}\textup{)} gives the exact microscopic count.
For large $\Delta$:
$\log d(\Delta) \sim 4\pi\sqrt{\Delta} -
\frac{3}{2}\log\Delta + \cdots$,
where the subleading term is the universal one-loop correction
\textup{(}Sen 2005, 2012\textup{)}.
The asymptotic expansion is the leading term of the Rademacher
exact formula for the Fourier coefficients of $1/\Phi_{10}$:
\[
d(\Delta) \;=\; \frac{2\pi}{\sqrt{\Delta}}
\sum_{c=1}^{\infty} \frac{K(c,\Delta)}{c}\,
I_{27/2}\!\Bigl(\frac{4\pi\sqrt{\Delta}}{c}\Bigr)
\;+\; \text{polar contributions},
\]
where $I_{27/2}$ is the modified Bessel function of the first
kind of order~$27/2$ and $K(c,\Delta)$ is a Kloosterman-type
sum over Sp$(4,\bZ)$. The $c = 1$ term dominates at large
$\Delta$, giving
$d(\Delta) \sim e^{4\pi\sqrt{\Delta}} \cdot \Delta^{-29/4}$.
The convergence of the Rademacher series is exponentially fast:
the $c = 2$ correction is suppressed by $e^{-2\pi\sqrt{\Delta}}$
relative to the leading term.
\end{proposition}

\begin{remark}[OSV conjecture]
\label{rem:osv}
\index{OSV conjecture}
The Ooguri--Strominger--Vafa relation
$Z_{\mathrm{BH}} = |Z_{\mathrm{top}}|^2$ connects the black hole
partition function to the topological string. For $K3 \times E$,
the topological string is exactly solvable at genus~$0$ and~$1$.
The shadow tower gives $F_g = 3 \cdot \lambda_g^{\mathrm{FP}}$
at each genus \textup{(all-weight, with cross-channel correction $\delta F_g^{\mathrm{cross}}$)}.
\end{remark}


%% ---------------------------------------------------------------
%% BRIDGE 2: Fourier--Jacobi as genus escalation
%% ---------------------------------------------------------------

\subsection{Fourier--Jacobi as genus escalation}
\label{subsec:bridge-fourier-jacobi}
\index{Fourier--Jacobi expansion!genus escalation bridge}

\begin{proposition}[Bridge: elliptic $R$-matrix data determines BKM
root system; \ClaimStatusProvedElsewhere]
\label{prop:bridge-fj-genus-escalation}
The Fourier--Jacobi expansion of the Igusa cusp form
$\Phi_{10}(\Omega)
= \sum_{m \geq 1} \phi_m(\tau,z)\,p^m$
uses Jacobi forms $\phi_m(\tau,z)$ on $E_\tau$. The leading
coefficient
\begin{equation}\label{eq:bridge-phi-1}
\phi_1(\tau,z) = \eta(\tau)^{18}\,\vartheta_1(\tau,z)^2
\end{equation}
is built from the same theta function $\vartheta_1(\tau,z)$ that defines the
elliptic propagator $d\log\vartheta_1(z|\tau)$
\textup{(}\eqref{eq:ell-propagator}\textup{)} and
Felder's elliptic $R$-matrix entries $a(u) = \theta(u{+}\eta)/\theta(\eta)$
\textup{(}Definition~\textup{\ref{def:elliptic-r})}.

The passage from genus-$1$ data \textup{(}Jacobi forms on $E_\tau$\textup{)}
to genus-$2$ data \textup{(}the Siegel form $\Phi_{10}$ on $\mathfrak{H}_2$\textup{)}
is the Borcherds multiplicative lift. Its input is the K3 elliptic genus
$Z_{K3} = 2\,\phi_{0,1}$
\textup{(}Definition~\textup{\ref{def:k3-elliptic-genus})}, and the exponents
$c_0(4nm - \ell^2)$ in the infinite product~\eqref{eq:borcherds-lift} are precisely
the root multiplicities of the BKM superalgebra $\mathfrak{g}_{\Delta_5}$
\textup{(}Proposition~\textup{\ref{prop:bkm-roots})}.

In the factorization $\eta^{18} = \eta^{24} \cdot \eta^{-6}$:
$\eta^{24} = \Delta$ is the Ramanujan discriminant
\textup{(}the $24$~free bosons of $A_E$, with $\kappa_{\mathrm{ch}}(A_E) = 24$\textup{)}, and
$\eta^{-6} = \eta^{-2\kappa_{\mathrm{ch}}}$ with $\kappa_{\mathrm{ch}} = \kappa_{\mathrm{ch}}(\Omega^{\mathrm{ch}}(K3 \times E))
= 3$. The shadow obstruction tower detects $\kappa_{\mathrm{ch}} = 3$
\textup{(}Proposition~\textup{\ref{prop:shadow-k3e})};
the Borcherds lift escalates this genus-$1$ datum to the genus-$2$ Siegel modular form.
\end{proposition}


%% ===============================================================
%% TWISTED HOLOGRAPHY
%% ===============================================================

\section{Twisted holography on \texorpdfstring{$K3 \times E$}{K3 x E}}
\label{sec:twisted-holography-k3xe}
\index{twisted holography!K3 x E|textbf}
\index{Costello--Li programme}

\begin{construction}[Costello--Li KS boundary algebra]
\label{constr:costello-li-k3xe}
\index{Kodaira--Spencer theory!K3 x E}
The Costello--Li programme extracts a boundary chiral algebra from
the HT twist of Type IIB on $K3 \times E$. Reduction of
Kodaira--Spencer theory along the topological directions (K3)
produces a chiral algebra $A_E$ on $E$ with generators from
harmonic forms on K3:
\begin{center}
\small
\renewcommand{\arraystretch}{1.2}
\begin{tabular}{lcl}
$h^{p,q}(K3)$ & Contribution & Generators \\
\hline
$h^{0,0} = 1$ & Scalar boson & $1$ \\
$h^{2,0} = 1$ & Symplectic boson & $1$ \\
$h^{1,1} = 20$ & Weight-$1$ bosons & $20$ \\
$h^{0,2} = 1$ & Conjugate boson & $1$ \\
$h^{2,2} = 1$ & Volume boson & $1$
\end{tabular}
\end{center}
Total: $c(A_E) = \chi(K3) = 24$. This is NOT the K3 sigma model
($c = 6$) but $24$ free bosons from the B-model reduction.
$\kappa_{\mathrm{ch}}(A_E) = 24$ \textup{(}rank of the free-boson lattice\textup{)},
$\kappa_{\mathrm{ch}}(A_E^!) = -24$ \textup{(}$\kappa_{\mathrm{ch}} + \kappa_{\mathrm{ch}}' = 0$ for
free fields\textup{)}.
\end{construction}

\begin{warning}[OPE level matrix for KS boundary bosons]
\label{warn:ope-level-mukai}
\index{OPE level!Mukai pairing distinction}
\index{Mukai pairing!OPE level}
The $24$ free bosons $J^a(z)$ ($a = 1, \ldots, 24$) of the KS
boundary algebra $A_E$ have OPE
$J^a(z)\,J^b(w) \sim \delta^{ab}/(z-w)^2$,
so the OPE level matrix is $\delta^{ab}$ (the unit matrix),
\emph{not} the Mukai pairing $G^{ab}$ of signature $(4,20)$ on
$H^*(K3, \Z)$. The Mukai pairing enters through the vertex
operators $V_\gamma(z) = {:}\!\exp(i\gamma_a J^a(z))\!{:}$
for lattice vectors $\gamma \in \Gamma_{\mathrm{Mukai}}$,
whose correlation functions involve
$\langle V_\gamma V_{\gamma'} \rangle \propto
\exp(\gamma \cdot G \cdot \gamma')$. Confusing the OPE level
$\delta^{ab}$ with the lattice pairing $G^{ab}$ produces wrong
$\kappa_{\mathrm{ch}}$ values: the OPE level gives $\kappa_{\mathrm{ch}}(A_E) = 24$
(sum of diagonal entries), while the Mukai norm
$\mathrm{tr}(G) = 4 - 20 = -16$ would give a wrong and
negative~$\kappa_{\mathrm{ch}}$.
\end{warning}

\begin{remark}[Holographic modular Koszul datum]
\label{rem:holo-koszul-k3xe}
\index{holographic modular Koszul datum!K3 x E}
$H(K3 \times E) = (A, A^!, C, r(z), \Theta_A,
\nabla^{\mathrm{hol}})$ with
$A = A_E$ ($c = 24$, $\kappa_{\mathrm{ch}} = 24$),
$A^! = A_E^!$ ($\kappa_{\mathrm{ch}} = -24$),
$C = Z^{\mathrm{der}}_{\mathrm{ch}}(A)$ (universal bulk),
$r(z) = \kappa_{\mathrm{ch}}\,\Omega/z$ (Casimir, $24$-dim: level prefix $\kappa_{\mathrm{ch}} = 24$),
$\Theta_A$ (bar-intrinsic MC element),
$\nabla^{\mathrm{hol}}$ (shadow connection, class~G).
\end{remark}

\begin{proposition}[Boundary-to-sigma ratio;
\ClaimStatusProvedHere]
\label{prop:boundary-sigma-ratio}
\index{modular characteristic!boundary vs sigma model}
\index{K3 x E!boundary-to-sigma ratio}
The boundary algebra $A_E$ and the sigma model VOA
$V_{K3}$ have modular characteristics in ratio~$12$.
At genus~$1$:
\begin{equation}\label{eq:boundary-sigma-ratio}
\frac{\kappa_{\mathrm{ch}}(A_E)}{\kappa_{\mathrm{ch}}(V_{K3})}
\;=\; \frac{24}{2} \;=\; 12,
\qquad
\frac{F_1(A_E)}{F_1(V_{K3})}
\;=\; 12.
\end{equation}
Here $A_E$ has $c = 24$ from the $24$ free bosons of the
Kaluza--Klein reduction on K3
\textup{(}Construction~\textup{\ref{constr:costello-li-k3xe})}
indexed by the Mukai lattice $H^*(K3, \Z)$, and
$V_{K3}$ is the $\cN = 4$ SCA at $c = 6$ with
$\kappa_{\mathrm{ch}}(V_{K3}) = 2$ \textup{(}Proposition~\textup{%
\ref{prop:kappa-k3})}. The SUSY Ward identities of the
$\cN = 4$ algebra reduce $\kappa_{\mathrm{ch}}$ by a factor of $2/3$
relative to the Virasoro formula $\kappa_{\mathrm{ch}}(\mathrm{Vir}_6) = 3$
\textup{(}Remark~\textup{\ref{rem:ap48-k3}}\textup{)}.

For $A_E$ \textup{(}$24$~free bosons, class~G,
uniform weight\textup{)}, the scalar formula
$F_g(A_E) = 24 \cdot \lambda_g^{\mathrm{FP}}$ holds at all
genera.
For $V_{K3}$ \textup{(}class~M, generators at weights
$1, 3/2, 2$\textup{)}, the scalar formula
$F_g = \kappa_{\mathrm{ch}} \cdot \lambda_g^{\mathrm{FP}}$ is proved at
genus~$1$; at $g \geq 2$ the multi-weight genus expansion
receives cross-channel corrections $\delta F_g^{\mathrm{cross}}$
\textup{(all-weight; Vol~I)}. The ratio~$12$ is
therefore proved at $g = 1$ and conditional at $g \geq 2$.
\end{proposition}

\begin{proof}
$\kappa_{\mathrm{ch}}(A_E) = 24$ (rank of the free-boson lattice).
$\kappa_{\mathrm{ch}}(V_{K3}) = 2$ (Proposition~\ref{prop:kappa-k3}).
Their ratio is $24/2 = 12$. At genus~$1$,
$F_1(\cA) = \kappa_{\mathrm{ch}}(\cA)/24$ for both algebras, giving the
ratio $12$. At $g \geq 2$: $A_E$ is uniform-weight (class~G),
so $F_g(A_E) = 24 \cdot \lambda_g^{\mathrm{FP}}$ at all genera
(UNIFORM-WEIGHT; class~G, scalar formula exact).
$V_{K3}$ has generators at weights $1, 3/2, 2$; the scalar
formula $F_g = 2 \cdot \lambda_g^{\mathrm{FP}}$ is proved at
$g = 1$ but receives multi-weight cross-channel corrections
$\delta F_g^{\mathrm{cross}}$ at $g \geq 2$ \textup{(all-weight; Vol~I)}.
\end{proof}

\begin{remark}[Leech connection]
\label{rem:leech-connection}
\index{Leech lattice!K3 connection}
The boundary partition function
$Z_1(A_E;\tau) = 1/\eta(\tau)^{48}$ at genus~$1$ matches the
Leech lattice VOA $V_\Lambda$ through order $q^1$:
$1/\eta^{48} = q^{-2}(1 + 48\,q + \cdots)$, while
$V_\Lambda$ has $\dim (V_\Lambda)_1 = 0$ and
$\dim (V_\Lambda)_2 = 196560$. The two diverge at
order~$q^2$: the $196560$ vectors of the Leech lattice have
no counterpart in the $24$-boson Fock space
at weight~$2$.
The Leech lattice has $\kappa_{\mathrm{ch}}(V_\Lambda) = 24 = \kappa_{\mathrm{ch}}(A_E)$
\textup{(}both are $24$ free bosons at level~$1$\textup{)},
but their weight-$2$ spaces differ.
\end{remark}


%% ===============================================================
%% MOCK MODULAR FORMS AND THE BKM SUPERALGEBRA
%% ===============================================================

\section{Mock modular forms and the BKM root system}
\label{sec:mock-modular-bkm}
\index{mock modular form!BPS|textbf}
\index{BKM superalgebra|textbf}

\subsection{BKM root multiplicities}
\label{subsec:bkm-roots}

\begin{proposition}[Root multiplicities; \ClaimStatusProvedElsewhere]
\label{prop:bkm-roots}
\index{BKM algebra!root multiplicities}
The BKM superalgebra $\mathfrak{g}_{\Delta_5}$ has root
multiplicities from the K3 elliptic genus:
\begin{itemize}
\item Real roots ($D = -1$, $\alpha^2 = 2$):
 $\mathrm{mult} = c_0(-1) = 2$.
\item Lightlike roots ($D = 0$, $\alpha^2 = 0$):
 $\mathrm{mult} = c_0(0) = 10$
 \textup{(}$20$ from $2\phi_{0,1}$\textup{)}.
\item Imaginary roots ($D > 0$):
 $\mathrm{mult} = c_0(D)$, unbounded.
 Leading: $c_0(3) = -64$, $c_0(4) = 108$,
 $c_0(7) = -513$, $c_0(8) = 808$.
\end{itemize}
Negative multiplicities at odd $D$ indicate fermionic
directions in the superalgebra. The discriminant coefficients
$c_0(D)$ of $\phi_{0,1}$ through $D = 12$ are
\textup{(}Eichler--Zagier convention\textup{)}:
\begin{center}
\small
\renewcommand{\arraystretch}{1.2}
\begin{tabular}{r|rrrrrrrrrrrrrr}
$D$ & $-1$ & $0$ & $3$ & $4$ & $7$ & $8$
 & $11$ & $12$ & $15$ & $16$ & $19$ & $20$ \\
\hline
$c_0(D)$ & $2$ & $10$ & $-64$ & $108$ & $-513$ & $808$
 & $-2752$ & $4016$ & $-11775$ & $16524$ & $-43263$ & $58956$
\end{tabular}
\end{center}
The $K3$ elliptic genus uses $2\,\phi_{0,1}$, so the BKM root
multiplicities are $\mathrm{mult}(\alpha) = 2\,c_0(D)$ for
discriminant $D$. The theta decomposition
$\phi_{0,1} = h_0(\tau)\,\vartheta_{0,1}(\tau,z)
+ h_1(\tau)\,\vartheta_{1,1}(\tau,z)$ separates even
($\ell$ even, $D \equiv 0 \pmod{4}$) and odd
($\ell$ odd, $D \equiv 3 \pmod{4}$) sectors, with
$h_0(\tau) = 10 + 108\,q + 808\,q^2 + \cdots$ and
$h_1(\tau) = q^{-1/4}(1 - 64\,q + \cdots)$.
\end{proposition}

\begin{proposition}[Root multiplicity constancy; \ClaimStatusProvedHere]
\label{prop:root-mult-constancy}
\index{BKM algebra!root multiplicity constancy}
Define the \emph{level} of a positive root $\alpha = (n, l, m)$ as
$\ell(\alpha) = n + m$. For all levels $\ell \geq 1$:
\[
 \sum_{\substack{\alpha \in \Delta_+ \\ \ell(\alpha) = \ell}}
 \mathrm{mult}(\alpha) \;=\; 24 \;=\; \chi(K3).
\]
\end{proposition}

\begin{proof}
At each level~$\ell$, the positive roots decompose into boundary
roots (those with $n = 0$ or $m = 0$) and interior roots (those
with $n, m \geq 1$).

\medskip\noindent\textbf{Boundary contribution.}
Roots with $m = 0$: $(n, l, 0)$ for $n = \ell$, varying~$l$.
The multiplicity is $c_0(-l^2)$, which is nonzero only for
$l = 0$ (giving $c_0(0) = 10$) and $l = \pm 1$ (giving
$c_0(-1) = 1$ each). Total: $10 + 1 + 1 = 12$.
Similarly for $n = 0$: total~$12$. Boundary total:~$24$.

\medskip\noindent\textbf{Interior cancellation.}
For each pair $(n, m)$ with $n + m = \ell$ and $n, m \geq 1$:
\[
 \sum_{l \in \bZ} c_0(4nm - l^2) \;=\; 0.
\]
This is the modular constancy of $\phi_{0,1}$ at $z = 0$.
The weak Jacobi form $\phi_{0,1}(\tau, z)$ has weight~$0$ and
index~$1$; evaluating at $z = 0$ gives a modular form of
weight~$0$ on $\mathrm{SL}_2(\bZ)$, which is necessarily constant
(the space $M_0(\mathrm{SL}_2(\bZ))$ is one-dimensional, spanned
by~$1$). The value is $\phi_{0,1}(\tau, 0) = 12$
(Eichler--Zagier, \emph{The Theory of Jacobi Forms}, 1985, p.~108).
At the level of Fourier coefficients:
$\sum_l c_0(4N - l^2) = 0$ for all $N \geq 1$, since the $q^N$
coefficient of the constant function~$12$ vanishes for $N \geq 1$.
The boundary gives~$24$, the interior gives~$0$, so the total at
every level is~$24$.
\end{proof}

\begin{remark}[The interior cancellation at level~$2$]
\label{rem:level-2-cancellation}
\index{BKM algebra!interior cancellation}
The interior at level~$2$ consists of roots $(1, l, 1)$ with
$\mathrm{mult} = c_0(4 - l^2)$. From the Eichler--Zagier table:
$l = 0$: $D = 4$, $c_0(4) = 108$;
$l = \pm 1$: $D = 3$, $c_0(3) = -64$ each;
$l = \pm 2$: $D = 0$, $c_0(0) = 10$ each.
Sum: $108 - 2 \cdot 64 + 2 \cdot 10 = 108 - 128 + 20 = 0$.
The cancellation is a direct consequence of
$\sum_l c_0(4 - l^2) = 0$ (the $q^1$ coefficient of
$\phi_{0,1}(\tau, 0) = 12$).
\end{remark}

\begin{remark}[Root multiplicity constancy and the shadow tower]
\label{rem:constancy-shadow}
\index{root multiplicity!shadow tower connection}
The constancy $\sum_{\ell(\alpha) = \ell} \mathrm{mult}(\alpha)
= 24 = \chi(K3)$ reflects the topological datum $\chi(K3)$
persisting through the BKM construction. This is precisely the
datum that the abelian $(2,0)$ pushforward
(Construction~\ref{constr:pushforward-chiral-algebra}) captures:
the Euler characteristic of the fiber enters through the lattice
rank and the partition function of the free-field theory. The
pushforward gives the correct topological backbone
($\chi = 24$) even though it gives the wrong shadow class
(class~$\mathbf{G}$ instead of class~$\mathbf{M}$) and the wrong
$\kappa_{\mathrm{ch}}$ ($\mathrm{rank}(\Lambda)$ instead of
$\dim_\bC(K3) = 2$).
\end{remark}


%% ---------------------------------------------------------------
%% BRIDGE 3: Hilbert schemes
%% ---------------------------------------------------------------

\subsection{The Hilbert scheme bridge}
\label{subsec:bridge-hilbert-schemes}
\index{Hilbert scheme!bridge to toroidal algebras}

\begin{theorem}[Toroidal action on $K$-theory of Hilbert schemes
{\cite{SV13}}; \ClaimStatusProvedElsewhere]
\label{thm:bridge-sv-toroidal}
\index{Schiffmann--Vasserot theorem}
\index{toroidal algebra!Hilbert scheme action}
Let $S$ be a smooth projective surface. The toroidal algebra
$U_{q,t}(\hat{\hat{\mathfrak{gl}}}_1)$ acts on
$\bigoplus_{n \geq 0} K({\operatorname{Hilb}}^n(S))$
by correspondences on the nested Hilbert scheme
$\operatorname{Hilb}^{n,n+1}(S)$ \textup{(}Schiffmann--Vasserot\textup{)}.
Specializations:
\begin{enumerate}[label=\textup{(\roman*)}]
\item \emph{Nakajima operators} \textup{(}$q = 1$\textup{)}: the
 Heisenberg algebra $\bigoplus_{n \in \bZ} \bQ\,\mathfrak{a}_n$
 acts on $\bigoplus_n H^*(\operatorname{Hilb}^n(S))$ via
 Nakajima correspondences,
 with $\mathfrak{a}_{-n}|0\rangle$
 the class of the punctual Hilbert scheme.
\item \emph{Haiman's theorem} \textup{(}$S = \bC^2$\textup{)}:
 the DAHA $\ddot{H}_n(q,t)$, a quotient of
 $U_{q,t}(\hat{\hat{\mathfrak{gl}}}_1)$, acts on
 $K(\operatorname{Hilb}^n(\bC^2))$ with fixed-point basis
 indexed by partitions $\lambda \vdash n$, and the
 character of the natural representation is the Macdonald
 polynomial $\widetilde{H}_\lambda(x;q,t)$.
\item \emph{$S = K3$}:
 $\chi(\operatorname{Hilb}^n(K3)) = p_{24}(n)$
 \textup{(}$24$-coloured partitions\textup{)}, so
 $\sum_{n \geq 0} \chi(\operatorname{Hilb}^n(K3))\,p^n
 = \prod_{m \geq 1} (1 - p^m)^{-24}$.
\end{enumerate}
\end{theorem}

\begin{proposition}[DMVV from Hilbert scheme characters;
\ClaimStatusProvedElsewhere]
\label{prop:bridge-dmvv-hilb}
\index{DMVV formula!Hilbert scheme}
The DVV formula~\eqref{eq:dvv}
$1/\Phi_{10} = \sum_N p^N\,\chi(\operatorname{Sym}^N(K3);\tau,z)$
is the generating function of equivariant Euler characteristics of
$\operatorname{Hilb}^N(K3)$ refined by the $\mathrm{SU}(2)_R$~fugacity
$z$ and the $L_0$~fugacity $q = e^{2\pi i\tau}$. The $24$~coloured
partitions arise from Nakajima's Heisenberg action
\textup{(}$24 = \chi(K3) = b_0 + b_2 + b_4 = 1 + 22 + 1$\textup{)}:
each colour corresponds to a cohomology class of K3, matching the
$24$~free bosons of the boundary algebra $A_E$
\textup{(}Construction~\textup{\ref{constr:costello-li-k3xe})}.
\end{proposition}

\begin{remark}[Synthesis: toroidal algebras meet BPS counting]
\label{rem:bridge-synthesis}
\index{Hilbert scheme!synthesis}
The three strands of this chapter converge on the Hilbert scheme.
The toroidal algebra $U_{q,t}(\hat{\hat{\mathfrak{g}}})$ of
\S\ref{sec:double-affine} provides the symmetry algebra
acting on $K(\operatorname{Hilb}^n(S))$.
The elliptic $R$-matrix of
\S\ref{sec:elliptic-quantum} appears as the transition
matrix between stable-envelope bases of $K(\operatorname{Hilb}^n(S))$
(Aganagic--Okounkov).
The BKM root multiplicities $c_0(D)$ of
Proposition~\ref{prop:bkm-roots} are the Fourier coefficients of the
K3~elliptic genus $Z_{K3} = 2\,\phi_{0,1}$, which in turn
counts representations of the Heisenberg subalgebra on
$\bigoplus_n H^*(\operatorname{Hilb}^n(K3))$.
The shadow invariant $\kappa_{\mathrm{ch}} = 3$ measures the single-copy
chiral algebra $\Omega^{\mathrm{ch}}(K3 \times E)$; the BKM algebra
$\mathfrak{g}_{\Delta_5}$ encodes the second-quantized
(symmetric-product) extension, and the gap between them is the content
of the shadow--Siegel gap analysis
(\S\ref{subsec:shadow-siegel-gap}).
\end{remark}


\subsection{Fourier--Jacobi expansion}
\label{subsec:fj-expansion}

\begin{remark}[Fourier--Jacobi structure]
\label{rem:fj-structure}
\index{Fourier--Jacobi expansion!Phi_{10}}
$\Phi_{10} = \sum_{m \geq 1} \phi_m(\tau,z)\,p^m$ with:
$\phi_1 = \phi_{10,1} = \eta^{18}\,\vartheta_1^2 =
-\Delta \cdot \phi_{-2,1}$
(unique Jacobi cusp form of weight~$10$, index~$1$);
$\phi_2 = 2\,\Delta\,\phi_{-2,1}\,\phi_{0,1}$
($\dim J_{10,2}^{\mathrm{cusp}} = 1$).
The factorization $\eta^{18} = \eta^{24} \cdot \eta^{-6}$
separates the Ramanujan discriminant from
$\eta^{-2\kappa_{\mathrm{ch}}}$ with $\kappa_{\mathrm{ch}} = 3$. The Jacobi form ring
$J_{*,*}^{\mathrm{weak}} = \bC[E_4, E_6, \phi_{-2,1},
\phi_{0,1}]$ (Eichler--Zagier).
\end{remark}


%% ===============================================================
%% GRAND SYNTHESIS AND OPEN PROBLEMS
%% ===============================================================

\section{Grand synthesis and open problems}
\label{sec:grand-synthesis-k3xe}
\index{K3 x E!grand synthesis}

\begin{remark}[Master data table for $K3 \times E$]
\label{rem:master-data-k3xe}
The complete invariant data assembled from the $39$ CY compute
engines:
\begin{center}
\small
\renewcommand{\arraystretch}{1.3}
\begin{tabular}{ll}
\textbf{Invariant} & \textbf{Value} \\
\hline
$\dim_\bC(K3 \times E)$ & $3$ \\
$\chi(K3 \times E)$ & $0$ \\
$h^{1,1} = h^{2,1}$ & $21$ \\
$\mathrm{rank}\,K_0$ & $48$ \\
$\dim \mathrm{HH}^*$ & $96$ \\
$\kappa_{\mathrm{ch}} = \kappa_{\mathrm{ch}}(\Omega^{\mathrm{ch}})$ & $3$ \\
$\kappa_{\mathrm{BKM}} = \mathrm{wt}(\Phi_{10})/2$ & $5$ \\
$\kappa_{\mathrm{BCOV}} = \chi/24$ & $0$ \\
Shadow depth & class M ($r_{\max} = \infty$) \\
$F_1$ & $1/8$ \\
$F_2$ & $7/1920$ \\
$F_3$ & $31/322560$ \\
$Z_{\mathrm{DT},0}$ & $1$ \\
$n^0_1$ (genus-$0$ BPS, $h = 1$) & $24$ \\
$S_{\mathrm{BH}}(\Delta)$ & $4\pi\sqrt{\Delta}$ \\
$c(A_E^{\mathrm{KS}})$ (twisted holography) & $24$ \\
$\kappa_{\mathrm{ch}}(A_E^{\mathrm{KS}})$ & $24$ \\
$\mathrm{Br}(K3 \times E)$ & $\supset (\bQ/\bZ)^{22}$
\end{tabular}
\end{center}
\end{remark}

\begin{definition}[$\kappa_{\mathrm{ch}}$-spectrum]
\label{def:kappa-spectrum}
\index{kappa-spectrum@$\kappa_{\mathrm{ch}}$-spectrum}
\index{modular characteristic!$\kappa_{\mathrm{ch}}$-spectrum}
For a Calabi--Yau threefold $X$, the \emph{$\kappa_{\mathrm{ch}}$-spectrum} is
\[
 \operatorname{Spec}_{\kappa_{\mathrm{ch}}}(X)
 \;:=\;
 \bigl\{\kappa_{\mathrm{ch}}(\cA) : \cA \text{ a chirally Koszul algebra
 associated to } X\bigr\}.
\]
Each element is the modular characteristic
(Vol~I, Theorem~D) of a specific algebraization of~$X$.
The $\kappa_{\mathrm{ch}}$-spectrum is an invariant of $X$ that
remembers the \emph{set} of modular characteristics arising
from all chirally Koszul algebraizations, not any single one.
\end{definition}

\begin{example}[$K3 \times E$: four-element $\kappa_{\mathrm{ch}}$-spectrum]
\label{ex:kappa-spectrum-k3xe}
\index{modular characteristic!polysemy}
The geometry $K3 \times E$ admits at least four distinct chirally
Koszul algebras, each arising from a different construction
:
\begin{center}
\small
\begin{tabular}{lccl}
 \toprule
 \textbf{Algebraization} & $c$ & $\kappa_{\mathrm{ch}}$ & \textbf{Shadow class} \\
 \midrule
 Chiral de Rham $\Omega^{\mathrm{ch}}(K3 \times E)$ & $0$ & $3$ & M \\
 K3 sigma model ($\cN = 4$ SCA) & $6$ & $2$ & M \\
 KS boundary algebra ($24$ free bosons) & $24$ & $24$ & G \\
 BPS spectrum ($\tfrac{1}{2}\operatorname{wt}(\Phi_{10})$)
 & --- & $5$ & M \\
 \bottomrule
\end{tabular}
\end{center}
Hence $\operatorname{Spec}_{\kappa_{\mathrm{ch}}}(K3 \times E) \supseteq \{2, 3, 5, 24\}$.
The four values probe different aspects of the geometry:
$\kappa_{\mathrm{ch}} = 3 = \dim_\bC(K3 \times E)$ is the chiral
dimension; $\kappa_{\mathrm{ch}}(\cA_{K3}) = 2$ is the K3 sigma model modular characteristic
	(an algebraization invariant, not a manifold invariant; cf.\);
$\kappa_{\mathrm{ch}}(A_E) = 24$ counts free bosons;
$\kappa_{\mathrm{BKM}} = 5$ is half the weight of the Igusa cusp form
$\Phi_{10}$.
\end{example}

\begin{remark}[The $\kappa_{\mathrm{ch}}$-spectrum as invariant]
\label{rem:kappa-spectrum-k3xe}
\index{kappa-spectrum}
The $\kappa_{\mathrm{ch}}$-spectrum (Definition~\ref{def:kappa-spectrum}) is
analogous to the multiple spectra of a Riemannian
manifold (Laplacian, Dirac, length): different constructions
applied to the same geometry probe different aspects of the modular
structure. Whether $\operatorname{Spec}_{\kappa_{\mathrm{ch}}}(X)$ is a complete
invariant (i.e., whether
$\operatorname{Spec}_{\kappa_{\mathrm{ch}}}(X) = \operatorname{Spec}_{\kappa_{\mathrm{ch}}}(X')$
implies $X \simeq X'$) is open.
\end{remark}

\begin{remark}[The second-quantization gap]
\label{rem:second-quantization-gap}
\index{second-quantization gap}
The shadow tower is first-quantized (single-copy).
$1/\Phi_{10}$ is second-quantized (DMVV symmetric product).
The ratio $\kappa_{\mathrm{BKM}}/\kappa_{\mathrm{ch}} = 5/3$
reflects the passage from single copy to $\mathrm{Sym}^N$
and is specific to $K3 \times E$ (the quintic gives a different
ratio). The Borcherds lift converts additive genus-$1$ data
into multiplicative genus-$2$ data: shadow genus escalation.
\end{remark}

\begin{proposition}[The BPS modular characteristic;
\ClaimStatusProvedHere]
\label{prop:kappa-bps-decomposition}
\index{modular characteristic!BPS}
\index{second quantization!kappa decomposition}
The BPS modular weight is
\begin{equation}\label{eq:kappa-bps-decomposition}
\kappa_{\mathrm{BKM}}
\;=\; \tfrac{1}{2}\,c_0(0)
\;=\; \tfrac{1}{2}\operatorname{wt}(\Phi_{10})
\;=\; 5
\end{equation}
\textup{(}Borcherds weight theorem;
Proposition~\textup{\ref{prop:bkm-weight-universal})}.
For $K3 \times E$ specifically, the numerical coincidence
$\kappa_{\mathrm{BKM}} = \kappa_{\mathrm{ch}}(K3) + \kappa_{\mathrm{ch}}(K3 \times E) = 2 + 3 = 5$
holds, where $\kappa_{\mathrm{ch}}(K3) = 2$ is the fiber modular
characteristic \textup{(}Proposition~\textup{\ref{prop:kappa-k3})}
and $\kappa_{\mathrm{ch}}(K3 \times E) = 3 = \dim_\bC$ is the
threefold modular characteristic
\textup{(}Proposition~\textup{\ref{prop:shadow-k3e})}.
This additive decomposition is NOT universal: it fails for
$\bZ/N\bZ$-orbifolds with $N \geq 2$
\textup{(}Remark~\textup{\ref{rem:bkm-decomposition-adversarial})}.
The ratio
\[
\frac{\kappa_{\mathrm{BKM}}}{\kappa_{\mathrm{ch}}}
\;=\; \frac{5}{3}
\;=\; \frac{\operatorname{wt}(\Phi_{10})}{2\,\dim_\bC(K3 \times E)}
\]
reflects the passage from single-copy to symmetric-product.
The distinction between $\operatorname{Hilb}^N$ and
$\operatorname{Sym}^N$ accounts for additional BPS states:
$\chi(\operatorname{Hilb}^2(K3)) - \chi(\operatorname{Sym}^2(K3))
= 324 - 300 = 24 = \chi(K3)$, arising from the exceptional
divisor of the Hilbert--Chow morphism
$\operatorname{Hilb}^2(K3) \to \operatorname{Sym}^2(K3)$.
The DMVV formula~\eqref{eq:dvv} uses the orbifold
$\operatorname{Sym}^N$ formula; the G\"ottsche formula gives
$\sum_N q^N \chi(\operatorname{Hilb}^N(K3))
= \prod_{k \geq 1}(1-q^k)^{-24}$.
\end{proposition}

\begin{proof}
The Borcherds weight theorem gives
$\kappa_{\mathrm{BKM}} = c_0(0)/2 = 10/2 = 5$ unconditionally.
The numerical coincidence
$\kappa_{\mathrm{ch}}(K3) + \kappa_{\mathrm{ch}}(K3 \times E) = 2 + 3 = 5$
is verified by independent computation of each summand
(Propositions~\ref{prop:kappa-k3} and~\ref{prop:shadow-k3e}).
That the sum agrees with $c_0(0)/2$ is specific to $K3 \times E$
($N = 1$ orbifold): the adversarial engine
\texttt{kappa\_bkm\_adversarial.py} confirms failure at $N \geq 2$.
The Hilbert--Chow difference: $\operatorname{Hilb}^2(S)$
for a K3 surface~$S$ is the blowup of $\operatorname{Sym}^2(S)$
along the diagonal, with exceptional divisor
$\bP(T_S) \cong \bP^1$-bundle. By the blowup formula,
$\chi(\operatorname{Hilb}^2) = \chi(\operatorname{Sym}^2)
+ \chi(S) = \binom{25}{2} + 24 = 300 + 24 = 324$.
\end{proof}

\begin{remark}[Hilbert--Chow difference at higher $N$]
\label{rem:hilbert-chow-higher}
\index{Hilbert--Chow!higher N}
The identity
$\chi(\operatorname{Hilb}^2(S)) - \chi(\operatorname{Sym}^2(S))
= \chi(S)$ generalizes: for all $N \geq 1$,
$\chi(\operatorname{Hilb}^N(S))$ is the coefficient of $q^N$ in
$\prod_{k \geq 1}(1-q^k)^{-\chi(S)}$ (G\"ottsche's formula),
while $\chi(\operatorname{Sym}^N(S))
= \binom{\chi(S) + N - 1}{N}$ (the combinatorial formula for
orbifold Euler characteristics). The difference
$\delta_N = \chi(\operatorname{Hilb}^N) -
\chi(\operatorname{Sym}^N)$ grows with~$N$; the leading
corrections come from the exceptional loci of the Hilbert--Chow
resolution $\operatorname{Hilb}^N(S) \to \operatorname{Sym}^N(S)$,
which become increasingly complicated as partitions of~$N$ acquire
more parts. For $K3$: $\delta_1 = 0$, $\delta_2 = 24 = \chi(K3)$,
$\delta_3 = 600$, $\delta_4 = 8100$, matching the G\"ottsche
coefficients $1, 24, 324, 3200, 25650, \ldots$ minus the
symmetric product coefficients
$1, 24, 300, 2600, 17550, \ldots$.
\end{remark}

\begin{openproblem}[Does the bar complex of $\cA_{K3}$ detect
$M_{24}$?]
\label{op:bar-m24}
\index{Mathieu moonshine!bar complex detection}
The mock modular half-multiplicities $a_n = 45, 231, 770,\ldots$
are dimensions of $M_{24}$ representations. The bar complex
$B(\cA_{K3})$ detects $\kappa_{\mathrm{ch}} = 2$ at degree~$2$ and the cubic
shadow at degree~$3$. Does the full shadow obstruction tower
$\Theta_{\cA_{K3}}$ encode the $M_{24}$ action? The mock modular
shadow $24\,\eta^3$ and the bar-complex shadow connection
$\nabla^{\mathrm{sh}}$ share the structural feature of encoding
modular non-invariance. Whether they are related by a precise
functor remains open.
\end{openproblem}

\begin{openproblem}[Factorization envelope of the K3 chiral algebra]
\label{op:fact-envelope-k3}
\index{factorization envelope!K3}
The factorization envelope $U^{\mathrm{mod}}_X(\cA_{K3})$ in the
sense of Nishinaka (2025/26) and Vicedo (2025) is not yet
constructed. At the genus-$0$ level, the K3 sigma model provides a
factorization algebra on $\mathrm{Ran}(X)$ for any algebraic curve
$X$; the modular completion (incorporating all genera) is the
frontier. The six-fold platonic package
$\Pi_X(\cA_{K3})$ is well-defined
(Construction~\ref{constr:platonic-package}); what remains is the
modular factorization envelope as the left adjoint of $\mathrm{Prim}^{\mathrm{mod}}$.
\end{openproblem}

\begin{remark}[Constructive reconstruction of $D^b(K3 \times E)$]
\label{rem:db-k3xe-reconstruction}
\index{derived category!K3 constructive reconstruction}
\index{Cech descent!K3}
The derived category $D^b(K3 \times E) \simeq D^b(K3) \otimes
D^b(E)$ is constructively reconstructible by three independent
methods:
\begin{enumerate}[label=\textup{(\roman*)}]
\item \emph{\v{C}ech descent.}
Two affine charts suffice for a smooth K3 surface
\textup{(}any smooth projective surface over an algebraically
closed field is covered by at most two affines\textup{)}.
The \v{C}ech descent spectral sequence
\textup{(}Theorem~\textup{\ref{thm:cech-descent})}
\[
E_1^{p,q} = \bigoplus_{|I|=p+1}
 \mathrm{HH}^q(D^b(U_I))
\;\Longrightarrow\;
\mathrm{HH}^{p+q}(D^b(K3))
\]
recovers $D^b(K3)$ from local data on $U_0, U_1$
with transition data on $U_{01}$.
\item \emph{BKR equivalence.}
At the Kummer point $K3 = T^4/\bZ_2$, the
Bridgeland--King--Reid theorem gives
$D^b(\operatorname{Hilb}^2(T^4/\bZ_2)) \simeq
D^b_{\bZ_2}(T^4)$, providing an explicit equivariant
description. The equivariant category
$D^b_{\bZ_2}(T^4)$ is generated by equivariant
line bundles and is computationally accessible.
\item \emph{Brauer obstruction.}
$\mathrm{Br}(K3) \supset (\bQ/\bZ)^{21}$
\textup{(}Remark~\textup{\ref{rem:brauer-k3xe})} obstructs
only \emph{twisted} derived categories. The
untwisted category $D^b(K3)$ is unobstructed.
\end{enumerate}
A structural constraint: CY threefolds carry no exceptional
objects \textup{(}$\chi(\cO_X, \cO_X) = 0$ by the
antisymmetry of Serre duality on a CY$_3$\textup{)}, so
$D^b(K3 \times E)$ admits no semiorthogonal decomposition
and must be treated as an indecomposable block
\textup{(}Proposition~\textup{\ref{prop:sod-k3xe})}.
\end{remark}

\begin{openproblem}[Global CY category from finitely many quiver charts]
\label{op:global-cy-category}
\index{Calabi--Yau category!global from quiver charts}
Given the constructive methods of
Remark~\ref{rem:db-k3xe-reconstruction}, can
$D^b(K3 \times E)$ be assembled from finitely many quiver
charts with explicit $A_\infty$-bimodule transition data?
The McKay correspondence provides the local charts at
orbifold points; the cluster mutation structure controls
transitions; the KS wall-crossing governs the DT invariant
bookkeeping. A constructive algorithm would connect the
local-to-global CY gluing of
\S\ref{sec:cy-local-global} to the global enumerative data of
\S\ref{sec:enum-geom-k3xe}.
\end{openproblem}


%% ---------------------------------------------------------------
\subsection{Research programmes}
\label{subsec:cy-research-programmes}
\index{K3 x E!research programmes}

The preceding computation identifies the shadow obstruction tower
of $K3 \times E$ down to precise numerical invariants. Ten research
programmes emerge, each probing a different direction of the
interface between bar-cobar duality and Calabi--Yau geometry.
These are recorded as formal conjectures and open problems;
each is grounded in computations from the $42$ CY engines
of the compute layer.

\bigskip

%%% --- Programme A: Constructive CY gluing ---

\begin{openproblem}[Programme A: constructive Calabi--Yau gluing]
\label{op:programme-a-cy-gluing}
\index{Calabi--Yau category!constructive gluing}
\index{quiver chart!constructive gluing}
The \v{C}ech descent of
\S\ref{sec:cy-local-global} reconstructs $D^b(K3 \times E)$
from local quiver charts in principle. Three questions sharpen
the programme.
\begin{enumerate}[label=\textup{(A\arabic*)}]
\item \emph{Minimum chart number.} For a K3 surface~$S$ of
 Picard rank~$\rho$, what is the minimum number of
 ADE-type quiver charts whose $A_\infty$-bimodule
 transition data recovers $D^b(S)$? The Mukai lattice
 $\widetilde{H}(S,\bZ) \cong U^4 \oplus E_8(-1)^2$ has
 rank~$24$; categorical generation requires a spanning class
 for thick subcategory generation, not merely $K$-theory
 spanning. For orbifold K3 (Kummer $T^4/\bZ_2$),
 the BKR equivalence
 $D^b(\operatorname{Kum}(A)) \simeq D^b_{\bZ/2}(A)$
 provides a single global chart; for smooth K3 without
 ADE singularities, no full exceptional collection exists
 (Bridgeland 1999), and one must generate via spherical
 objects and their twist functors.
\item \emph{Brauer obstruction.} The gluing obstruction
 lives in $H^2(S, \cO_S) \cong \bC$
 (the Brauer class). When does a nontrivial Brauer
 class kill the descent? For generic K3 with $\rho = 0$,
 $\operatorname{Br}(S)$ contains $(\bQ/\bZ)^{22}$;
 for $K3 \times E$,
 $\operatorname{Br}(K3 \times E)
 \supset (\bQ/\bZ)^{22+1}$.
\item \emph{Beyond ADE.} The McKay correspondence
 handles orbifold singularities. For smooth K3 at a
 generic moduli point (no singularities), the relevant
 local charts are formal neighborhoods of subvarieties,
 not ADE resolution charts. Does the spherical twist
 functor orbit of $\cO_S$ generate $D^b(S)$ with
 finitely many transitions?
\end{enumerate}
\end{openproblem}

\begin{conjecture}[ADE chart decomposition of K3]
\label{conj:ade-chart-k3}
\index{K3 surface!ADE chart decomposition}
\ClaimStatusConjectured\quad
Every algebraic K3 surface admits a finite ADE-chart
decomposition of its derived category: there exist finitely
many open sets $\{U_\alpha\}$ with tilting generators
$T_\alpha$ on each $U_\alpha$, and
$A_\infty$-bimodule transition data on overlaps, such that
$D^b(S) \simeq \operatorname{holim}_{\check{C}} D^b(U_\alpha)$.
The minimum number of charts is bounded below by
$\lceil \operatorname{rank}\,K_0(S) / \max_\alpha
\operatorname{rank}\,K_0(\operatorname{End}(T_\alpha)) \rceil$.
\end{conjecture}


%%% --- Programme B: kappa as universal moonshine multiplier ---

\begin{openproblem}[Programme B: $\kappa_{\mathrm{ch}}$ as universal moonshine multiplier]
\label{op:programme-b-moonshine}
\index{moonshine!universal multiplier}
\index{modular characteristic!moonshine}
The Mathieu moonshine decomposition gives
$A_n = 2 \cdot \dim(\rho_n)$ with
$\rho_n \in \operatorname{Irr}(M_{24})$, where the
factor $2 = \kappa_{\mathrm{ch}}(\cA_{K3})$ is the modular characteristic
of the K3 sigma model
(Remark~\ref{rem:factor-2-is-kappa}).
This suggests a universal principle.
\begin{enumerate}[label=\textup{(B\arabic*)}]
\item \emph{Monster moonshine.} For the Monster module
 $V^\natural$ ($c = 24$, $\kappa_{\mathrm{ch}}(V^\natural) = 12$ by
 Remark~\ref{rem:lattice-voa-k3} and
 Vol~I, Theorem~\ref{thm:lattice-voa-bar}, applied to the
 Leech lattice construction):
 does $A_n^{\text{Monster}} = \kappa_{\mathrm{ch}}(V^\natural) \cdot
 \dim(\rho_n^{\mathbb{M}})$ for the Monster group
 $\mathbb{M}$?
 The McKay--Thompson series $T_g(\tau) =
 \sum_n \operatorname{tr}(\rho_n(g))\,q^{n-1}$ is a
 genus-zero Hauptmodul; the question is whether the
 \emph{dimensions} (not characters) factor through~$\kappa_{\mathrm{ch}}$.
\item \emph{Conway moonshine.} For $V^{s\natural}$
 ($c = 12$, the shorter moonshine module):
 $\kappa_{\mathrm{ch}}(V^{s\natural}) = ?$ and does
 $A_n = \kappa_{\mathrm{ch}} \cdot \dim(\rho_n^{Co_0})$?
 The Conway group $Co_0$ acts on the Leech lattice;
 $V^{s\natural}$ is the $\bZ/2$-orbifold of the
 rank-$12$ lattice VOA.
\item \emph{General principle.} For a holomorphic VOA~$V$
 of central charge~$c$ with finite automorphism group~$G$,
 does the elliptic genus decompose as
 $Z_g(\tau,z) = \kappa_{\mathrm{ch}}(V) \cdot \sum_n
 \dim(\rho_n)\,\operatorname{ch}_n(g)$ for $g \in G$?
\end{enumerate}
\end{openproblem}

\begin{conjecture}[Universal moonshine multiplier]
\label{conj:universal-moonshine-multiplier}
\index{moonshine!universal multiplier conjecture}
\ClaimStatusConjectured\quad
For a holomorphic VOA~$V$ with finite automorphism group~$G$,
the twined partition function decomposes as
\[
Z_g(\tau,z)
\;=\;
\kappa_{\mathrm{ch}}(V) \cdot \sum_{n \geq 0}
\dim(\rho_n)\,\mathrm{ch}_n(g),
\qquad g \in G,
\]
where $\kappa_{\mathrm{ch}}(V)$ is the bar-complex modular characteristic,
$\{\rho_n\}$ are virtual $G$-representations, and
$\mathrm{ch}_n$ are character functions determined by the
elliptic genus of~$V$. For $V = \cA_{K3}$ ($c = 6$) with
$G = M_{24}$: $\kappa_{\mathrm{ch}} = 2$ and
$A_n = 2 \cdot \dim(\rho_n)$ as established by Gannon~\cite{Gannon16}.
\end{conjecture}


%%% --- Programme C: Second-quantization bridge ---

\begin{openproblem}[Programme C: second-quantization bridge]
\label{op:programme-c-second-quantization}
\index{second quantization!shadow bridge}
\index{plethystic exponential}
The ratio $\kappa_{\mathrm{BKM}}/\kappa_{\mathrm{ch}} = 5/3$
(Proposition~\ref{prop:kappa-bps-decomposition}) encodes the
passage from first to second quantization for $K3 \times E$.
\begin{enumerate}[label=\textup{(C\arabic*)}]
\item \emph{Universality.} The additive decomposition
 $\kappa_{\mathrm{BKM}}
 = \kappa_{\mathrm{ch}}(\text{surface}) + \kappa_{\mathrm{ch}}(\text{CY}_3)$
 is a numerical coincidence for $K3 \times E$ ($N = 1$)
 and FAILS for $\bZ/N\bZ$-orbifolds with $N \geq 2$
 \textup{(}Remark~\textup{\ref{rem:bkm-decomposition-adversarial}};

 The correct universal formula is
 $\kappa_{\mathrm{BKM}} = c_N(0)/2$
 \textup{(}Borcherds weight theorem\textup{)}.
 For the quintic: $\kappa_{\mathrm{BKM}}$ is undefined
 \textup{(}no K3 fibration\textup{)}; the replacement invariant
 $\kappa_{\mathrm{BCOV}} = \chi/24 = -200/24$ must be extracted from
 the genus-$2$ topological string amplitude.
\item \emph{Plethystic shadow.} The DMVV formula
 $1/\Phi_{10} = \mathrm{PE}[2\,\phi_{0,1}]$ expresses
 the second-quantized partition function as the plethystic
 exponential of the single-copy elliptic genus. What is
 the shadow tower of $\mathrm{PE}[\cA]$ in terms of
 the shadow tower of~$\cA$? The tensor product gives
 $\kappa_{\mathrm{ch}}(\cA^{\otimes N}) = N\,\kappa_{\mathrm{ch}}(\cA)$ by additivity,
 but the $\bZ_N$-orbifold projection to $\mathrm{Sym}^N$
 modifies the tower through twisted sectors.
\item \emph{Adams operations.} The Adams operations
 $\psi^k$ act on the shadow partition function by
 rescaling:
 \begin{equation}\label{eq:adams-shadow}
 \psi^k\bigl(Z^{\mathrm{sh}}\bigr)(\hbar)
 \;=\; Z^{\mathrm{sh}}(k\hbar)
 \;=\; \exp\!\Bigl(\sum_{g \geq 1}
 k^{2g}\,F_g\,\hbar^{2g}\Bigr).
 \end{equation}
 This is the $\lambda$-ring structure on the shadow
 partition function: $\psi^k$ multiplies the genus-$g$
 amplitude by $k^{2g}$, consistent with the
 interpretation of $\hbar^{2g}$ as a genus-counting
 weight. The Adams operations are ring homomorphisms
 ($\psi^k \circ \psi^l = \psi^{kl}$) and commute
 with the $\hat{A}$-genus generating function
 (Vol~I, Theorem~\ref{thm:shadow-connection}).
 Does this intertwine with the Hecke operators $V_k$
 on the Jacobi form $\phi_{0,1}$ via the DMVV
 formula~\eqref{eq:dvv}? The Hecke transform
 $V_k(\phi_{0,1})$ and the Adams-rescaled shadow
 $\psi^k(Z^{\mathrm{sh}})$ would need to be related
 through the Borcherds product for the intertwining
 to be exact.
\end{enumerate}
\end{openproblem}

\begin{remark}[BPS modular characteristic: universal formula vs.\ coincidence]
\label{rem:kappa-bps-universality-vol3}
\index{modular characteristic!BPS universality}
The correct universal formula for K3-fibered CY$_3$s is
\begin{equation}\label{eq:kappa-bkm-universal-k3e}
\kappa_{\mathrm{BKM}}
\;=\;
\tfrac{1}{2}\,c_N(0)
\end{equation}
\textup{(}Borcherds weight theorem;
Proposition~\textup{\ref{prop:bkm-weight-universal})},
where $c_N(0)$ is the constant term of the $\bZ/N\bZ$-orbifold
elliptic genus. The additive decomposition
$\kappa_{\mathrm{BKM}} = \kappa_{\mathrm{ch}}(S) + \kappa_{\mathrm{ch}}(S \times E) = 2 + 3 = 5$
is a numerical coincidence for $N = 1$ \textup{(}unorbifolded $K3 \times E$\textup{)};
it FAILS for all $\bZ/N\bZ$-orbifolds with $N \geq 2$
\textup{(}Remark~\textup{\ref{rem:bkm-decomposition-adversarial}};
\texttt{kappa\_bkm\_adversarial.py}, $62$~tests\textup{)}.
For non-K3-fibered CY$_3$s \textup{(}quintic, $\bC^3$,
conifold, local $\bP^2$\textup{)}, $\kappa_{\mathrm{BKM}}$ is
undefined; replacement invariants are $\kappa_{\mathrm{BCOV}} = \chi(X)/24$
\textup{(}compact\textup{)} or shadow depth
\textup{(}all, conditional on CY-A\textup{)}.
\end{remark}


%%% --- Programme D: Schottky shadow programme ---

\begin{openproblem}[Programme D: Schottky shadow programme]
\label{op:programme-d-schottky}
\index{Schottky problem!shadow tower}
\index{shadow obstruction tower!Schottky}
The shadow tower lives on $\overline{\cM}_g$
(via tautological classes), while Siegel modular forms
live on $\cA_g$. For $g \leq 3$, the Torelli map
$j \colon \cM_g \hookrightarrow \cA_g$ is an
isomorphism ($g = 1$), birational ($g = 2$), or
injective with codimension-$0$ complement ($g = 3$).
For $g \geq 4$, the Jacobian locus $\cJ_g = j(\cM_g)$
has codimension $(g-2)(g-3)/2$ in $\cA_g$; the
Schottky problem is nontrivial. The explicit numerical
values are:
\begin{center}
\small
\renewcommand{\arraystretch}{1.2}
\begin{tabular}{ccccl}
$g$ & $\dim \cM_g$ & $\dim \cA_g$
 & $\mathrm{codim}(\cJ_g)$ & Status \\
\hline
$2$ & $3$ & $3$ & $0$ & Torelli birational \\
$3$ & $6$ & $6$ & $0$ & Torelli dominant \\
$4$ & $9$ & $10$ & $1$ & Schottky--Igusa $J$ \\
$5$ & $12$ & $15$ & $3$ & \\
$6$ & $15$ & $21$ & $6$ & \\
$7$ & $18$ & $28$ & $10$ & \\
$8$ & $21$ & $36$ & $15$ & \\
$9$ & $24$ & $45$ & $21$ & \\
$10$ & $27$ & $55$ & $28$ & Crossover: $\mathrm{codim} > \dim \cM_{10}$
\end{tabular}
\end{center}
The crossover at $g = 10$ (where
$\mathrm{codim}(\cJ_{10}) = 28 > 27 = \dim \cM_{10}$) marks
the genus beyond which the Jacobian locus is ``more absent
than present'' in~$\cA_g$. The shadow tower, living on
$\overline{\cM}_g$, sees a progressively thinner slice of
the Siegel modular world as~$g$ grows.
\begin{enumerate}[label=\textup{(D\arabic*)}]
\item \emph{Tautological insulation.} The shadow amplitude
 $F_g = \kappa_{\mathrm{ch}}\,\lambda_g^{\mathrm{FP}}$ \textup{(all-weight, with cross-channel correction $\delta F_g^{\mathrm{cross}}$)}
 lives in the
 tautological ring $R^*(\overline{\cM}_g)$, which is
 insensitive to the Schottky constraint. The shadow
 tower cannot detect whether a period matrix is a
 Jacobian: it is tautologically insulated from Schottky.
\item \emph{Physical partition function.}
 The physical genus-$g$ partition function
 $Z_g(\Omega)$ for $K3 \times E$ \emph{is} constrained
 by Schottky at $g \geq 4$: it is defined on $\cJ_g$,
 not all of $\cA_g$. The gap between shadow
 (tautological) and partition function (analytic)
 at $g \geq 4$ measures the failure of the shadow tower
 to capture the full moduli dependence.
\item \emph{The Schottky--Igusa form.} The Schottky--Igusa
 form $J \in S_8(\operatorname{Sp}(8,\bZ))$ vanishes
 exactly on $\cJ_4$ at genus~$4$. Its weight $8 = 2g$
 is significant. Does any shadow invariant at genus~$4$
 detect~$J$? The tautological ring at genus~$4$ has
 dimension $\dim R^4(\overline{\cM}_4) = 2$; the shadow
 contributes to one component, but $J$ lives transverse
 to the tautological ring.
\end{enumerate}
\end{openproblem}

\begin{conjecture}[Shadow tower generates the tautological
projection]
\label{conj:shadow-taut-projection}
\index{tautological ring!shadow generation}
\ClaimStatusConjectured\quad
At genus~$g$, the shadow obstruction tower of a class-M
algebra \textup{(}shadow depth $r_{\max} = \infty$\textup{)}
generates the image of the restriction map
$\operatorname{res} \colon S_*(\operatorname{Sp}(2g,\bZ))
\to R^*(\overline{\cM}_g)$
from the ring of Siegel modular forms to the tautological
ring. The shadow tower does not see the transverse
directions to $\cJ_g$ in $\cA_g$, but it generates the
full tautological projection of the Siegel data.
\end{conjecture}


%%% --- Programme E: Mock modularity and the bar complex ---

\begin{openproblem}[Programme E: mock modularity and the bar complex]
\label{op:programme-e-mock-modularity}
\index{mock modular form!bar complex}
\index{Appell--Lerch sum}
The mock modular form $h(\tau) =
2\,q^{-1/8}(-1 + 45q + 231q^2 + \cdots)$ encoding the
massive $\cN = 4$ multiplicities has shadow
$S(\tau) = 24\,\eta(\tau)^3$ (weight~$3/2$, cusp form).
\begin{enumerate}[label=\textup{(E\arabic*)}]
\item \emph{The factor $24$.} The shadow
 $S(\tau) = 24\,\eta^3$ contains
 $24 = \chi(K3)$, the topological Euler characteristic.
 Where does this factor arise in the bar complex?
 The bar curvature $\kappa_{\mathrm{ch}} = 2$ does not directly
 produce~$24$; the factor $24/\kappa_{\mathrm{ch}} = 12$ is the
 modular discriminant level. An algebraic derivation
 from the bar complex would connect the mock modular
 shadow to the bar-complex shadow connection
 $\nabla^{\mathrm{sh}}$.
\item \emph{The Appell--Lerch sum.} The massless $\cN = 4$
 character is controlled by the Appell--Lerch sum
 $\mu(\tau,z)$. The non-holomorphic completion
 $\hat{h}(\tau,\bar\tau)$ is modular of weight~$1/2$.
 Is there a bar-complex interpretation of~$\mu$?
 The Appell--Lerch sum regularizes the divergent
 theta-type sum; the bar complex regularizes the
 divergent Feynman-type sum.
 Both are regularizations of the same modular anomaly.
\item \emph{Half-integral weight.} The mock modular
 form~$h$ has weight~$1/2$; the shadow has weight~$3/2$.
 The bar-complex shadow connection $\nabla^{\mathrm{sh}}$
 produces integer-weight objects (the shadow metric $Q_L$
 has weight~$2$; the Faber--Pandharipande
 $\lambda_g$ are tautological, hence of integral weight).
 The passage to half-integral weight requires either
 a metaplectic extension or a theta correspondence.
 Concretely: is there a metaplectic bar complex whose
 shadow connection produces weight-$3/2$ objects?
\end{enumerate}
\end{openproblem}

\begin{warning}[Notation: $\mu$ is overloaded]
\label{warn:mu-overloaded}
\index{mu function!notation overload}
The symbol~$\mu$ denotes two distinct objects in the K3 literature:
\begin{enumerate}[label=\textup{(\roman*)}]
\item The \emph{Zwegers $\mu$-function}
 $\mu(\tau, z) = \sum_{n \in \Z}
 (-1)^n\,q^{n(n+1)/2}\,y^n/(1 - y\,q^n)$, the Appell--Lerch sum
 controlling the massless $\cN = 4$ character
 (Programme~E above).
\item The full \emph{$\cN = 4$ massless character}
 $\mathrm{ch}^{\mathrm{massless}}(\tau, z)
 = \mu(\tau, z) \cdot \vartheta_1(\tau, z)/\eta(\tau)^3$
 (with additional theta and eta factors).
\end{enumerate}
In the compute engines,
\texttt{cy\_mock\_modular\_bps\_engine.py} uses~$\mu$ for~(i)
while \texttt{cy\_n4sca\_k3\_engine.py} uses~$\mu$ for~(ii).
Care is needed when combining results from different engines.
The Zwegers $\mu$-function is \emph{not} a Jacobi form: it fails
the elliptic transformation law (it has a pole at
$z \in \Z + \Z\tau$). Its non-holomorphic completion
$\hat{\mu}(\tau, \bar\tau, z)$ is modular.
\end{warning}

\begin{conjecture}[Mock shadow tower]
\label{conj:mock-shadow-tower}
\index{mock modular form!mock shadow tower}
\ClaimStatusConjectured\quad
The mock modular completion
$\hat{h}(\tau,\bar\tau) = h(\tau) + h^*(\tau,\bar\tau)$
is the genus-$1$ projection of a ``mock shadow tower''
that combines bar-complex data with the Appell--Lerch
regularization. There exists an extension
$\widetilde{\Theta}_{\cA_{K3}}$ of the shadow
obstruction tower $\Theta_{\cA_{K3}}$ valued in mock
modular forms of weight~$1/2$, whose non-holomorphic
completion is modular and whose holomorphic part restricts
to $\kappa_{\mathrm{ch}}(\cA_{K3}) \cdot h(\tau)$ at genus~$1$.
The shadow of this mock extension is
$\kappa_{\mathrm{ch}} \cdot \chi(K3) \cdot \eta^3$,
where the product $\kappa_{\mathrm{ch}} \cdot \chi(K3) = 2 \cdot 24 = 48$
is the rank of the charge lattice
$\Gamma = H^{\mathrm{even}}(K3 \times E, \bZ)$.
\end{conjecture}


%%% --- Programme F: Modular factorization envelope ---

\begin{openproblem}[Programme F: modular factorization envelope]
\label{op:programme-f-factorization-envelope}
\index{factorization envelope!modular}
The factorization envelope $U_E(\cA_E) =
\mathrm{Sym}^{\mathrm{ch}}(\mathfrak{h}_{24})$ of the
KS boundary algebra produces $24$ free bosons on~$E$,
with character $1/\eta(\tau)^{24}$. The modular
completion $U_E^{\mathrm{mod}}(\cA_E)$ must incorporate
higher-genus data; its character at genus~$1$ should
receive corrections from the shadow tower.
\begin{enumerate}[label=\textup{(F\arabic*)}]
\item \emph{The Leech divergence.} The Leech lattice
 VOA $V_{\mathrm{Leech}}$ has character
 $J(\tau) + 24 = q^{-1}(1 + 196884\,q + \cdots)$.
 The free-field envelope $1/\eta^{24}
 = q^{-1}(1 + 24\,q + 324\,q^2 + \cdots)$ matches
 the Leech VOA character at $q^{-1}$ and $q^0$
 (after adding~$24$) but diverges at $q^1$:
 $324 \neq 196884$. The discrepancy is the interacting
 structure of the Leech lattice VOA beyond free fields.
\item \emph{Modular factorization envelope.} Is there a
 universal construction
 $U_X^{\mathrm{mod}}(L)$ for all cyclically admissible
 Lie conformal algebras~$L$
 (Vol~I, Definition~\ref{def:cyclically-admissible}) that
 carries a natural shadow obstruction tower? Nishinaka's
 construction (2025/26) handles the genus-$0$ level;
 the modular completion requires incorporating the
 $F_g$ data at all genera.
\item \emph{Super extension.} The $\cN = 4$ SCA at
 $c = 6$ is a super vertex algebra. Nishinaka's
 construction applies to ordinary Lie conformal algebras;
 the super extension (with fermionic generators $G^\pm$,
 $\widetilde{G}^\pm$) requires the super-factorization
 framework of Costello--Gwilliam. The factorization
 envelope of the super Lie conformal algebra underlying
 the $\cN = 4$ SCA is not yet constructed.
\end{enumerate}
\end{openproblem}

\begin{conjecture}[Modular factorization envelope existence]
\label{conj:modular-fact-envelope}
\index{factorization envelope!modular existence}
\ClaimStatusConjectured\quad
For every cyclically admissible Lie conformal algebra~$L$,
the modular factorization envelope
$U_X^{\mathrm{mod}}(L)$ exists as a factorization algebra
on $\mathrm{Ran}(X)$ equipped with the shadow obstruction
tower $\Theta_{U_X^{\mathrm{mod}}(L)}$. This envelope
is the left adjoint of the modular primitive-current functor
$\mathrm{Prim}^{\mathrm{mod}}$
\textup{(}Construction~\textup{\ref{constr:platonic-package})},
specialized to the CY setting.
\end{conjecture}


%%% --- Programme G: BKM algebras and scattering diagrams ---

\begin{openproblem}[Programme G: BKM algebras and scattering diagrams]
\label{op:programme-g-bkm-scattering}
\index{Borcherds--Kac--Moody algebra!scattering diagram}
\index{scattering diagram!BKM}
\index{Kontsevich--Soibelman!wall-crossing}
The BKM root system of
$K3 \times E$ lives in $\Gamma^{2,2} = U \oplus U$
(lattice of signature $(2,2)$, rank~$4$).
Root multiplicities are determined by $c_0(D)$ from the K3
elliptic genus. The Kontsevich--Soibelman scattering
diagram for $K3 \times E$ associates a wall to each BPS ray
with attached function determined by the root multiplicity.
\begin{enumerate}[label=\textup{(G\arabic*)}]
\item \emph{Two MC elements.} The shadow obstruction tower
 $\Theta_{\cA}$ is an MC element in the convolution
 algebra $\mathfrak{g}_\cA^{\mathrm{mod}}$; the KS
 automorphism $\mathfrak{S}$ is an MC element in the
 motivic Hall algebra $\cH(\cC)$. These are MC elements
 in different Lie algebras. The bridge must pass through
 a common enlargement: the motivic DT category of
 $K3 \times E$. At the naive level, the BCH
 pair-commutator does \emph{not} reproduce the
 $\phi_{0,1}$ multiplicities (quantitative
 mismatch by non-uniform ratios).
\item \emph{Motivic level.} The identification
 ``scattering diagram $=$ shadow obstruction tower'' holds at
 the motivic Hall algebra level (Kontsevich--Soibelman 2008),
 where the product on walls is the motivic
 quantum torus product. Instantiating this at the
 numerical/Euler characteristic level loses the
 higher-genus information that the motivic structure
 retains.
\item \emph{Tropical skeleton.} The tropical limit of
 the scattering diagram is a piecewise-linear fan in
 $\Gamma^{2,2} \otimes \bR$. This fan should be the
 tropical skeleton of the shadow obstruction tower
 restricted to BPS walls. The identification would
 connect Mok's log FM tropicalization
 (Theorem~\ref{thm:planted-forest-tropicalization})
 to the BPS scattering diagram.
\end{enumerate}
\end{openproblem}

\begin{conjecture}[Scattering diagram as tropical shadow]
\label{conj:scattering-tropical-shadow}
\index{scattering diagram!tropical shadow}
\ClaimStatusConjectured\quad
The scattering diagram of $K3 \times E$ is the tropical
skeleton of the shadow obstruction tower:
$\operatorname{Scat}(K3 \times E)
= \operatorname{Trop}(\Theta_{\cA_{K3 \times E}})$,
where the left side is the Kontsevich--Soibelman
scattering diagram \textup{(}with wall-crossing factors
from the BPS spectrum\textup{)} and the right side is the
tropicalization of the shadow obstruction tower restricted
to the charge lattice~$\Gamma^{2,2}$.
The identification holds at the level of the motivic Hall
algebra; at the numerical level, higher BPS bound-state
contributions must be included.
\end{conjecture}


%%% --- Descent foundations ---

\begin{remark}[Descent sufficiency and K-theory verification]
\label{rem:descent-sufficiency}
\index{descent theorem!Zariski sufficiency}
\index{K-theory!descent verification}
Three levels of descent apply to the CY category programme:
\begin{enumerate}[label=\textup{(\roman*)}]
\item \emph{Zariski suffices.} For smooth separated schemes~$X$,
 the natural functor
 $D^b(X) \to \operatorname{holim}_{\mathrm{\check{C}ech}}
 D^b(U_\alpha)$ is an equivalence
 (To\"en 2007, Lurie HA~2017, Theorem~19.4.5):
 coherent sheaves satisfy effective descent for the Zariski
 topology. No \'etale or flat refinement is needed.
\item \emph{ADE cocycles close.} For each ADE singularity type,
 the bimodule cocycle condition
 $M_{01} \otimes^L_B M_{12} \simeq M_{02}$ closes:
 \begin{center}
 \small
 \renewcommand{\arraystretch}{1.2}
 \begin{tabular}{lccl}
 Type & $|\Gamma|$ & $\dim\Pi_0(Q)$ & Cocycle mechanism \\
 \hline
 $A_n$ & $n{+}1$ & $\binom{n{+}2}{3}$ & Flop involution \\
 $D_n$ & $4(n{-}2)$ & Table~\ref{prop:preprojective-mckay}
 & Tilting generator $T = \bigoplus P_i$ \\
 $E_6$ & $24$ & $168$ & $\operatorname{End}(T) = \Pi(\check{Q})$ \\
 $E_7$ & $48$ & $384$ & $\operatorname{End}(T) = \Pi(\check{Q})$ \\
 $E_8$ & $120$ & $1080$ & $\operatorname{End}(T) = \Pi(\check{Q})$
 \end{tabular}
 \end{center}
 The mechanism is uniform: the tilting generator
 $T = \bigoplus_{i \in Q_0} P_i$ of the resolution
 $\widetilde{\bC^2/\Gamma}$ has endomorphism algebra
 $\operatorname{End}(T) \cong \Pi_0(Q)$, and the Morita
 equivalence $D^b(\Pi_0(Q)\text{-mod}) \simeq
 D^b(\widetilde{\bC^2/\Gamma})$ automatically closes the
 bimodule cocycle for any number of overlapping charts
 (Bridgeland--King--Reid for type~$A$;
 Kapranov--Vasserot, Van den Bergh for all types).
\item \emph{K-theory verification.} For $X = K3 \times E$,
 successful descent must recover
 $\operatorname{rk} K_0(K3 \times E) = 48 =
 \operatorname{rk} K_0(K3) \cdot \operatorname{rk} K_0(E)
 = 24 \cdot 2$. For each ADE chart of type~$\Gamma$:
 $\operatorname{rk} K_0(\mathrm{Jac}(Q_\Gamma, W_\Gamma))
 = |Q_0| = r + 1$ (the number of vertices in the McKay
 quiver). The Mayer--Vietoris sequence in K-theory
 provides the gluing verification
 ($117$ tests in \texttt{cy\_descent\_theorem\_engine.py}).
\end{enumerate}
\end{remark}


%%% --- Programme H: Descent theorem programme ---

\begin{openproblem}[Programme H: descent theorem programme]
\label{op:programme-h-descent}
\index{descent theorem!CY programme}
\index{A-infinity bimodule!descent}
The descent theorem for dg categories (To\"en 2007, Lurie HA)
guarantees that Zariski descent recovers $D^b(X)$ for smooth
separated schemes (Remark~\ref{rem:descent-sufficiency}).
The quiver chart descent of
\S\ref{sec:cy-local-global} uses $A_\infty$-bimodule transition
data on overlaps.
\begin{enumerate}[label=\textup{(H\arabic*)}]
\item \emph{Bimodule cocycle closure.} For three overlapping
 charts $U_0, U_1, U_2$ with transition bimodules
 $M_{01}, M_{12}, M_{02}$, the triple cocycle condition
 $M_{01} \otimes^L_B M_{12} \simeq M_{02}$ must close up
 to coherent homotopy. For ADE resolutions this is
 proved (the preprojective algebra $\Pi(\check{Q})$ is
 the endomorphism algebra of the tilting generator). For
 smooth K3 patches glued along a divisor, the bimodule
 cocycle is the restriction of the Atiyah class.
\item \emph{Obstruction.} The obstruction to extending a
 two-chart descent to a global descent lives in
 $\mathrm{HH}^2$ of the transition bimodule:
 $\operatorname{ob} \in
 \mathrm{HH}^2(M_{01}, M_{01})$. For K3, $h^{0,2} = 1$
 provides a one-dimensional obstruction space. When is
 the obstruction trivial?
\item \emph{Higher-dimensional CY.} For a smooth projective
 CY $d$-fold~$X$ with ADE singularities, does the quiver
 chart descent reconstruct $D^b(X)$? The $d = 2$ (K3)
 and $d = 3$ ($K3 \times E$) cases are established by
 the compute layer. The general case requires understanding
 the higher Hochschild obstructions.
\end{enumerate}
\end{openproblem}


%%% --- Programme I: Higher-dimensional shadows ---

\begin{openproblem}[Programme I: higher-dimensional CY shadows]
\label{op:programme-i-higher-dim}
\index{Calabi--Yau!higher-dimensional shadows}
\index{shadow obstruction tower!higher dimension}
The shadow tower of $K3 \times E$ ($\dim_\bC = 3$) connects
to genus-$2$ Siegel modular forms via the Borcherds lift.
What happens in higher CY dimension?
\begin{enumerate}[label=\textup{(I\arabic*)}]
\item \emph{$\mathrm{CY}_4 = K3 \times K3$.} By additivity,
 $\kappa_{\mathrm{ch}} = \kappa_{\mathrm{ch}}(K3) + \kappa_{\mathrm{ch}}(K3) = 2 + 2 = 4$.
 Shadow class: M (product of two class-M factors).
 The genus-$2$ amplitude $F_2 = 4 \cdot 7/5760 = 7/1440$.
 Is there a genus-$3$ Siegel modular form analogue of
 $\Phi_{10}$?
\item \emph{$\mathrm{CY}_5 = K3 \times K3 \times E$.}
 Similarly,
 $\kappa_{\mathrm{ch}} = 2 + 2 + 1 = 5$. The same
 $\kappa_{\mathrm{BKM}} = 5$
 as for $K3 \times E$ arises, but for different reasons:
 the CY$_5$ has $\kappa_{\mathrm{ch}} = 5$ directly, without
 second quantization. Does this numerical coincidence
 have algebraic content?
\item \emph{The dimensional ladder.} For a product
 $\mathrm{CY}_d = X_1 \times \cdots \times X_k$,
 additivity gives
 $\kappa_{\mathrm{ch}} = \sum_i \kappa_{\mathrm{ch}}(X_i) = \sum_i \dim_\bC(X_i) = d$.
 The shadow class should be
 $\max(\text{classes of } X_i)$ under the ordering
 $G < L < C < M$. Does this determine the shadow
 tower of the product completely?
\end{enumerate}
\end{openproblem}

\begin{conjecture}[CY product shadow decomposition]
\label{conj:cy-product-shadow}
\index{Calabi--Yau!product shadow}
\ClaimStatusConjectured\quad
For a product $\mathrm{CY}_d = X_1 \times \cdots \times X_k$ of
Calabi--Yau manifolds,
\begin{align}
\kappa_{\mathrm{ch}}(X_1 \times \cdots \times X_k)
 &\;=\; \sum_{i=1}^k \kappa_{\mathrm{ch}}(X_i)
 \;=\; \sum_{i=1}^k \dim_\bC(X_i)
 \;=\; d, \\
\operatorname{class}(X_1 \times \cdots \times X_k)
 &\;=\; \max_i \operatorname{class}(X_i)
\end{align}
\textup{(}under the ordering
$G < L < C < M$\textup{)}. The shadow obstruction tower
$\Theta_{X_1 \times \cdots \times X_k}$ at the scalar
level is determined by $\kappa_{\mathrm{ch}} = d$ alone; the higher-degree
components are controlled by the most complex factor.
\end{conjecture}


%%% --- Programme J: Convergence vs divergence ---

\begin{openproblem}[Programme J: convergence vs.\ divergence
and non-perturbative completions]
\label{op:programme-j-convergence}
\index{shadow partition function!convergence}
\index{topological string!divergence}
\index{resurgence}
The shadow partition function $Z^{\mathrm{sh}}(\hbar) =
\sum_{g \geq 1} F_g\,\hbar^{2g}$ converges absolutely for
$|\hbar| < 2\pi$ (Proposition~\ref{prop:shadow-gf-convergence}),
is entire after Borel summation, and shows no resurgent
structure. The convergence radius $2\pi$ arises because
$F_g = \kappa_{\mathrm{ch}} \cdot \lambda_g^{\mathrm{FP}}$ with
$\lambda_g^{\mathrm{FP}} \sim (2\pi)^{-2g}$, so
$Z^{\mathrm{sh}}$ converges where $|\hbar/(2\pi)| < 1$.
The function $(\hbar/2)/\sin(\hbar/2)$ has poles at
$\hbar = 2n\pi$ for $n \in \Z \setminus \{0\}$; this is
$\sin(\hbar/2)$, not $\sin(\hbar)$, so the radius is
$2\pi$, not~$\pi$.
The topological string partition function
$Z_{\mathrm{top}}(\hbar)$ diverges as $(2g)!$ (Gevrey-$1$),
is resurgent, and exhibits Stokes phenomena.
\begin{enumerate}[label=\textup{(J\arabic*)}]
\item \emph{The shadow as convergent piece.}
 The shadow tower extracts the \emph{tautological}
 content of the genus expansion: intersection numbers on
 $\overline{\cM}_g$. The topological string adds
 \emph{non-tautological} contributions (worldsheet
 instantons, non-perturbative effects). Is
 $Z_{\mathrm{top}} = Z^{\mathrm{sh}} \cdot (1 +
 Z_{\mathrm{np}})$ where $Z_{\mathrm{np}}$ encodes
 worldsheet instantons?
\item \emph{Non-perturbative corrections.} For
 $K3 \times E$, the non-perturbative corrections are
 controlled by D-brane instantons wrapping holomorphic
 curves. The GV invariants $n^g_\beta$ count the
 BPS content. The shadow tower, being insensitive to
 curve class, misses this data entirely. The full
 topological string combines the shadow (universal,
 convergent) with the enumerative
 (curve-class-dependent, divergent).
\item \emph{The OSV bridge.}
 The OSV conjecture $Z_{\mathrm{BH}} = |Z_{\mathrm{top}}|^2$
 relates the black hole partition function (computed via
 $1/\Phi_{10}$ for $K3 \times E$) to the topological
 string. Since $Z^{\mathrm{sh}}$ extracts the
 tautological piece of $Z_{\mathrm{top}}$, the shadow
 controls the tautological piece of the black hole entropy.
 The precise relationship at subleading order ($\log\Delta$
 corrections to $S_{\mathrm{BH}}$) is governed by
 $F_1 = \kappa_{\mathrm{ch}}/24$ (Sen's logarithmic correction
 formula).
\end{enumerate}
\end{openproblem}


\begin{proposition}[Class G algebras: shadow $=$ topological string]
\label{prop:class-g-no-instantons}
\index{class G!no instanton corrections}
\index{shadow depth!class G exactness}
\ClaimStatusProvedHere\quad
For class~G algebras \textup{(}shadow depth $r_{\max} = 2$,
the shadow tower terminates at degree~$2$\textup{)},
\begin{equation}\label{eq:class-g-exact}
F_g^{\mathrm{sh}}(\cA) \;=\; F_g^{\mathrm{top}}(\cA)
\qquad \text{for all } g \geq 1.
\end{equation}
The shadow tower \emph{is} the topological string for
free-field-type algebras: the shadow partition function
$Z^{\mathrm{sh}} = \exp\!\bigl(\sum_{g \geq 1}
F_g\,\hbar^{2g}\bigr)$ with
$F_g = \kappa_{\mathrm{ch}} \cdot \lambda_g^{\mathrm{FP}}$ receives no
worldsheet instanton corrections because the bar complex is
formal \textup{(}$m_n = 0$ for $n \geq 3$\textup{)},
so the generating function is the $\hat{A}$-genus
$(\hbar/2)/\sin(\hbar/2)$ evaluated at $\kappa_{\mathrm{ch}}$. For
class~L, C, and M algebras, the higher-degree shadow
components $S_3, S_4, \ldots$ contribute additional
``instanton-like'' corrections to $F_g$ at genus $g \geq 2$
that modify the free-field $\hat{A}$-genus formula.
For the KS boundary algebra $A_E$ of $K3 \times E$:
$A_E$ is $24$ free bosons, hence class~G, hence
$F_g^{\mathrm{sh}}(A_E)$ is exact. The K3 sigma model
$\cA_{K3}$ is class~M \textup{(}infinite depth\textup{)},
so its shadow tower receives corrections at all degrees.
\end{proposition}

\begin{proof}
For a class-G algebra, $\Theta_\cA^{\leq r} = \Theta_\cA^{\leq 2}$
for all $r \geq 2$ (the tower terminates). The shadow obstruction
tower collapses to its scalar projection $\kappa_{\mathrm{ch}}$: the genus-$g$
amplitude is $F_g = \kappa_{\mathrm{ch}} \cdot \lambda_g^{\mathrm{FP}}$ with no higher-degree corrections. This is the full content of the
bar-intrinsic MC element $\Theta_\cA$ for class G. Since the
topological string at the free-field point receives no worldsheet
instanton corrections (the moduli space of maps is trivial),
$F_g^{\mathrm{top}} = F_g^{\mathrm{sh}}$.
\end{proof}


\begin{proposition}[Bar functor commutes with homotopy colimits]
\label{prop:bar-hocolim}
\index{bar functor!homotopy colimit}
\index{homotopy colimit!bar complex}
\ClaimStatusProvedHere\quad
Let $\{A_\sigma\}_{\sigma \in I}$ be a diagram of dg algebras
indexed by a small category~$I$. Then
\begin{equation}\label{eq:bar-hocolim-k3xe}
B\!\left(\operatorname{hocolim}_{\sigma \in I} A_\sigma\right)
\;\simeq\;
\operatorname{hocolim}_{\sigma \in I}\, B(A_\sigma)
\end{equation}
as dg coalgebras \textup{(}quasi-isomorphism\textup{)}.
\end{proposition}

\begin{proof}
The bar functor $B$ is a left Quillen functor from augmented dg
algebras to conilpotent dg coalgebras
(Vol~I, Theorem~\ref{thm:bar-cobar-adjunction}). Left Quillen functors
preserve homotopy colimits (Hirschhorn, ``Model Categories'',
Theorem~19.4.5). The quasi-isomorphism~\eqref{eq:bar-hocolim-k3xe}
is the derived functor applied to the homotopy colimit diagram.
Computational verification: $145$ tests in
\texttt{bar\_hocolim\_chain\_level.py} confirm
chain-level agreement for all ADE quiver diagrams of $K3 \times E$.
\end{proof}

\begin{remark}[Descent for bar complexes]
\label{rem:bar-descent}
\index{bar complex!descent}
Proposition~\ref{prop:bar-hocolim} gives a descent theorem for
bar complexes: the global bar complex of the hocolim is recovered
from the local bar complexes and their transition data. For the
constructive gluing of $K3 \times E$ from quiver charts
(\S\ref{sec:cy-local-global}), this means the global shadow
obstruction tower $\Theta_{\cA_{K3 \times E}}$ can be assembled
from local shadow towers $\Theta_{\cA_\sigma}$ on each chart,
with transition maps controlled by the $A_\infty$-bimodule data
of Construction~\ref{constr:ainfty-bimodule}.
\end{remark}


\begin{remark}[Summary of research programmes]
\label{rem:cy-programme-summary}
\index{K3 x E!programme summary}
Programmes A--J span the interface between modular Koszul
duality and Calabi--Yau geometry. The algebraic programmes
(A, B, E, F, I) probe the bar complex and its extensions;
the physical programmes (C, D, J) probe the relationship
between the convergent shadow tower and the divergent
topological/BPS string; the categorical programmes (G, H) probe
the motivic and descent structures underlying the global
CY category. Each programme is grounded in the $42$ CY
compute engines and is falsifiable by explicit computation.
Cross-references to this chapter and to the foundations in Vols~I--II:
\begin{itemize}
\item Programme A: \S\ref{sec:cy-local-global},
 Open Problem~\ref{op:global-cy-category}.
\item Programme B: Open Problem~\ref{op:bar-m24},
 Proposition~\ref{prop:kappa-k3}.
\item Programme C: Proposition~\ref{prop:kappa-bps-decomposition},
 Remark~\ref{rem:second-quantization-gap}.
\item Programme D: Theorem~\ref{thm:shadow-siegel-gap},
 Vol~I, Theorem~\ref{thm:riccati-algebraicity}.
\item Programme E: Open Problem~\ref{op:bar-m24},
 Vol~I, Theorem~\ref{thm:shadow-connection}.
\item Programme F: Open Problem~\ref{op:fact-envelope-k3},
 Vol~I, Construction~\ref{constr:platonic-package}.
\item Programme G: \S\ref{sec:bkm-k3e},
 Vol~I, Theorem~\ref{thm:planted-forest-tropicalization}.
\item Programme H: Theorem~\ref{thm:cech-descent},
 \S\ref{sec:cy-local-global}.
\item Programme I: Vol~I, Theorem~D
 (Theorem~\ref{thm:genus-universality}),
 Vol~I, Definition~\ref{def:shadow-depth-classification}.
\item Programme J: Proposition~\ref{prop:shadow-gf-convergence},
 Vol~I, Theorem~\ref{thm:general-hs-sewing}.
\end{itemize}
\end{remark}


%% ---------------------------------------------------------------
\section{The Borcherds--Kac--Moody algebra of \texorpdfstring{$K3 \times E$}{K3 x E}}
\label{sec:bkm-k3e}
\index{Borcherds--Kac--Moody algebra|textbf}
\index{Igusa cusp form|textbf}
\index{K3 surface!BKM algebra}

The preceding sections developed the elliptic bar complex on a
single elliptic curve $E_\tau$. The following sections examine a fundamentally
different role for genus-$2$ modular forms: the Calabi--Yau
threefold $K3 \times E$ produces a Borcherds--Kac--Moody (BKM)
superalgebra whose denominator function is the Igusa cusp form
$\Phi_{10}$, a Siegel modular form of weight~$10$ on
$\operatorname{Sp}(4, \bZ)$. The shadow obstruction tower
of $K3 \times E$ detects the topological skeleton of this
structure, but the passage from shadow amplitudes to the full
BPS partition function encounters four independent obstructions.
The bridge, and its limitations, require precise quantification.

\subsection{The Igusa cusp form and the DVV formula}
\label{subsec:igusa-dvv}

\begin{definition}[Igusa cusp form]
\label{def:igusa-cusp-form}
\index{Igusa cusp form!definition}
The \emph{Igusa cusp form} $\Phi_{10}$ is the unique Siegel
cusp form of weight~$10$ on $\operatorname{Sp}(4,\bZ)$:
\[
\dim S_{10}\bigl(\operatorname{Sp}(4,\bZ)\bigr) = 1.
\]
It admits three equivalent constructions:
\begin{enumerate}[label=\textup{(\roman*)}]
\item \emph{Product of even theta characteristics.}
$\Phi_{10}(\Omega) = -2^{-12} \prod_{\text{even}\,m}
\vartheta[m](\Omega)$, where the product runs over the~$10$
even theta characteristics of genus~$2$.
\item \emph{Fourier--Jacobi expansion.}
$\Phi_{10}(\Omega) = \sum_{m \geq 1}
\phi_m(\tau,z)\,p^m$ with $p = e^{2\pi i \sigma}$ and
$\Omega = \bigl(\begin{smallmatrix} \tau & z \\ z & \sigma
\end{smallmatrix}\bigr) \in \mathfrak{H}_2$.
The leading coefficient is
\begin{equation}\label{eq:phi-10-1}
\phi_1 \;=\; \phi_{10,1}(\tau,z)
\;=\; \eta(\tau)^{18}\,\vartheta_1(\tau,z)^2,
\end{equation}
the unique Jacobi cusp form of weight~$10$, index~$1$.
\item \emph{Borcherds multiplicative lift.}
Let $Z_{K3}(\tau,z) = 2\,\phi_{0,1}(\tau,z)$ be the
$K3$ elliptic genus, where $\phi_{0,1}$ is the unique
weak Jacobi form of weight~$0$, index~$1$ (Eichler--Zagier
normalization: $\phi_{0,1}(\tau,0) = 12$). Then
\begin{equation}\label{eq:borcherds-lift}
\Phi_{10}(\Omega) \;=\; p\,q\,y \prod_{\substack{
(n,\ell,m) > 0}}
\bigl(1 - q^n\,y^\ell\,p^m\bigr)^{c_0(4nm - \ell^2)},
\end{equation}
where $q = e^{2\pi i \tau}$, $y = e^{2\pi i z}$,
$p = e^{2\pi i \sigma}$, and $c_0(D)$ are the Fourier
coefficients of the $K3$ elliptic genus indexed by
discriminant~$D$.
\end{enumerate}
\end{definition}

\begin{remark}[Conventions]\label{rem:bkm-conventions}
\index{conventions!BKM}
We follow the Eichler--Zagier convention throughout.
In the DVV convention, $\Phi_{10}$ may differ
by a sign; the product formula~\eqref{eq:borcherds-lift}
and the uniqueness of $S_{10}(\operatorname{Sp}(4,\bZ))$
fix the normalization up to scalar. The elliptic genus
$Z_{K3} = 2\,\phi_{0,1}$ evaluates to
$Z_{K3}(\tau,0) = 24 = \chi(K3)$ at $z = 0$.
\end{remark}

The DVV formula (Dijkgraaf--Verlinde--Verlinde,
\cite{DMVV}) identifies the $\frac{1}{4}$-BPS
partition function of Type~IIA on $K3 \times T^2$:
\begin{equation}\label{eq:dvv}
Z_{\frac{1}{4}\text{-BPS}}(\Omega) \;=\; \frac{1}{\Phi_{10}(\Omega)}
\;=\; \sum_{N \geq 0}\,p^N\,\chi\bigl(\operatorname{Sym}^N(K3);
\tau,z\bigr),
\end{equation}
where the second equality is the DMVV infinite-product formula
expressing $1/\Phi_{10}$ as the generating function for the
orbifold elliptic genera of the symmetric products
$\operatorname{Sym}^N(K3)$.

\begin{proposition}[Genus-$2$ nature of $\Phi_{10}$;
\ClaimStatusProvedElsewhere]
\label{prop:phi10-genus2}
\index{Igusa cusp form!genus-2}
The Igusa cusp form $\Phi_{10}$ is \emph{intrinsically}
a genus-$2$ object: it is a Siegel modular form on
$\operatorname{Sp}(4,\bZ)$, whose natural domain is the
Siegel upper half-space
$\mathfrak{H}_2 = \{\Omega \in \mathrm{Mat}_{2 \times 2}(\bC)
: \Omega = \Omega^t,\; \operatorname{Im}(\Omega) > 0\}$
of complex dimension~$3$. The three
variables $(\tau,z,\sigma)$ of the period matrix
$\Omega = \bigl(\begin{smallmatrix} \tau & z \\
z & \sigma\end{smallmatrix}\bigr)$ encode the three
chemical potentials conjugate to the electric, magnetic,
and dyonic charge invariants of $\frac{1}{4}$-BPS states.
No genus-$1$ modular form captures the full dyonic
spectrum: the Fourier--Jacobi coefficient $\phi_{10,1}$
is the $m = 1$ term in the $p$-expansion, retaining only
the single-centered contribution.
\end{proposition}

\subsection{The Borcherds multiplicative lift}
\label{subsec:borcherds-lift}

\begin{theorem}[Borcherds lift; \ClaimStatusProvedElsewhere]
\label{thm:borcherds-lift-k3e}
\index{Borcherds lift|textbf}
\index{multiplicative lift}
Let $\phi_{0,1}(\tau,z)$ be the unique weak Jacobi form
of weight~$0$, index~$1$. The Borcherds multiplicative
lift
\[
\operatorname{Borch}\colon J_{0,1}^{\mathrm{weak}}
\;\longrightarrow\;
S_{10}\bigl(\operatorname{Sp}(4,\bZ)\bigr)
\]
sends $2\,\phi_{0,1}$ to $\Phi_{10}$, producing the
infinite product~\textup{\eqref{eq:borcherds-lift}}.
This is the genus-$1$ to genus-$2$ bridge: the
genus-$1$ datum $Z_{K3} = 2\,\phi_{0,1}$ encodes
the exponents of the genus-$2$ Siegel modular form
$\Phi_{10}$.
\end{theorem}

\begin{remark}[Structure of the lift]
\label{rem:borcherds-lift-structure}
The Borcherds lift converts additive data (Fourier
coefficients of $Z_{K3}$) into multiplicative data
(exponents in the product formula for $\Phi_{10}$).
The exponent $c_0(D)$ at discriminant $D = 4nm - \ell^2$
is the root multiplicity of the BKM superalgebra
$\mathfrak{g}_{\Delta_5}$: real roots ($D < 0$) have
multiplicity~$c_0(-1) = 2$, lightlike roots ($D = 0$)
have multiplicity $c_0(0) = 10$ (or $20$ from the full
$K3$ elliptic genus $2\,\phi_{0,1}$), and imaginary roots
($D > 0$) have unbounded multiplicities.

The leading Fourier--Jacobi coefficient
$\phi_1 = \phi_{10,1} = \eta^{18}\,\vartheta_1^2$
(equation~\eqref{eq:phi-10-1}) is itself a product
of genus-$1$ objects. The factorization
$\eta^{18} = \eta^{24} \cdot \eta^{-6}$ separates
the Ramanujan discriminant $\Delta = \eta^{24}$
(controlling the modular discriminant) from
$\eta^{-6} = \eta^{-2\kappa_{\mathrm{ch}}}$ with $\kappa_{\mathrm{ch}} = 3$
(the shadow tower's genus-$1$ partition function
$Z_1(\tau) = \eta(\tau)^{-2\kappa_{\mathrm{ch}}}$ for
$K3 \times E$ at $\kappa_{\mathrm{ch}} = \dim_\bC(K3 \times E) = 3$).
\end{remark}

\subsection{Shadow tower data for \texorpdfstring{$K3 \times E$}{K3 x E}}
\label{subsec:shadow-tower-k3e}

The modular Koszul framework assigns to $K3 \times E$
a well-defined shadow obstruction tower. We record
the explicit values.

\begin{proposition}[Shadow amplitudes of $K3 \times E$;
\ClaimStatusProvedHere]
\label{prop:shadow-k3e}
\index{K3 surface!shadow amplitudes}
Let $\cA = \Omega^{\mathrm{ch}}(K3 \times E)$ be the chiral
de~Rham complex of $K3 \times E$. The modular characteristic
is $\kappa_{\mathrm{ch}}(\cA) = \dim_\bC(K3 \times E) = 3$
\textup{(}this is not $c/2$ in general; here the central charge is
$c = 2 \cdot 3 = 6$, and $c/2 = 3$ coincides with $\kappa_{\mathrm{ch}}$
only because $\dim_\bC = 3$\textup{)}.
The shadow amplitudes \textup{(}Vol~I, Theorem~D;
Theorem~\textup{\ref{thm:modular-characteristic})}, in the
generating function convention $\sum_{g \geq 1} F_g\,\hbar^{2g}$,
are
\begin{equation}\label{eq:shadow-k3e}
F_g(K3 \times E) \;=\; \kappa_{\mathrm{ch}} \cdot \lambda_g^{\mathrm{FP}}
\;=\; 3 \cdot \frac{(2^{2g-1}-1)\,|B_{2g}|}{2^{2g-1}\,(2g)!}
\quad \textup{(all-weight, with cross-channel correction $\delta F_g^{\mathrm{cross}}$ at $g \geq 2$)},
\end{equation}
where $\lambda_g^{\mathrm{FP}}$ is the Faber--Pandharipande
intersection number. Explicitly:
\begin{center}
\small
\renewcommand{\arraystretch}{1.2}
\begin{tabular}{cccccc}
$g$ & $1$ & $2$ & $3$ & $4$ & $5$ \\
\hline
$\lambda_g^{\mathrm{FP}}$ & $\frac{1}{24}$ & $\frac{7}{5760}$
 & $\frac{31}{967680}$ & $\frac{127}{154828800}$
 & $\frac{73}{3503554560}$ \\[4pt]
$F_g$ & $\frac{1}{8}$ & $\frac{7}{1920}$
 & $\frac{31}{322560}$ & $\frac{127}{51609600}$
 & $\frac{73}{1167851520}$
\end{tabular}
\end{center}
The genus-$1$ amplitude $F_1 = \kappa_{\mathrm{ch}}/24 = 1/8$ matches the
leading behaviour of $Z_1(\tau) = \eta(\tau)^{-6}$, since
$\log\eta(\tau) = 2\pi i \tau / 24 + O(q)$ and the
$\eta$-power is $-2\kappa_{\mathrm{ch}} = -6$.
\end{proposition}

\begin{proof}
Direct from Vol~I, Theorem~D
(Theorem~\ref{thm:modular-characteristic})
with $\kappa_{\mathrm{ch}} = 3$. The value $\kappa_{\mathrm{ch}} = 3$ is the modular
characteristic of $\Omega^{\mathrm{ch}}(K3 \times E)$,
equal to $\dim_\bC(K3 \times E)$ by the Chern character
computation for the chiral de~Rham complex of a Calabi--Yau
manifold. Each $F_g$ is verified by three independent
paths: direct Bernoulli computation, the $\hat{A}$-generating
function $(t/2)/\sin(t/2) - 1 = \sum_{g \geq 1}
\lambda_g^{\mathrm{FP}}\,t^{2g}$, and the additivity
$\kappa_{\mathrm{ch}}(K3 \times E) = \kappa_{\mathrm{ch}}(K3) + \kappa_{\mathrm{ch}}(E) = 2 + 1 = 3$
(Vol~I, Corollary~\ref{cor:shadow-extraction}).
\end{proof}

\begin{proposition}[Convergence of the shadow generating function;
\ClaimStatusProvedHere]
\label{prop:shadow-gf-convergence}
\index{shadow generating function!convergence radius}
\index{Borel summability!shadow tower}
The shadow generating function
$Z^{\mathrm{sh}}(t) := \kappa_{\mathrm{ch}}\bigl((t/2)/\sin(t/2) - 1\bigr)
= \sum_{g \geq 1} F_g\,t^{2g}$ has the following analytic
properties:
\begin{enumerate}[label=\textup{(\roman*)}]
\item \emph{Convergence radius.}
$R = 2\pi$. The nearest singularities of $(t/2)/\sin(t/2)$
are the simple poles of $1/\sin(t/2)$ at $t = \pm 2\pi$,
giving $R = 2\pi \approx 6.283$.
\item \emph{Closed form.}
$Z^{\mathrm{sh}}(t) = \kappa_{\mathrm{ch}}\bigl((t/2)/\sin(t/2) - 1\bigr)$,
an algebraic function of~$\sin(t/2)$.
\item \emph{Borel transform.}
The Borel transform $\hat{Z}(u)
= \sum_{g \geq 1} F_g\,u^{2g}/(2g)!$ is an
\emph{entire} function of~$u$ \textup{(}no singularities
on $\bR_+$\textup{)}. There is no resurgence: the
Borel sum recovers $Z^{\mathrm{sh}}$ exactly.
\end{enumerate}
This contrasts sharply with the topological string partition
function $Z_{\mathrm{top}}(\lambda) = \exp\bigl(\sum_{g \geq 0}
F_g^{\mathrm{top}}\,\lambda^{2g-2}\bigr)$, which is
Gevrey-$1$ divergent: $F_g^{\mathrm{top}} \sim (2g)!$ from
the volume growth of $\mathcal{M}_g$.
For $K3 \times E$ with $\kappa_{\mathrm{ch}} = 3$:
$Z^{\mathrm{sh}}(t) = 3\bigl((t/2)/\sin(t/2) - 1\bigr)$
converges absolutely for $|t| < 2\pi$.
\end{proposition}

\begin{proof}
The function $f(t) = (t/2)/\sin(t/2)$ is meromorphic with poles
at $t = 2\pi k$ for $k \in \bZ \setminus \{0\}$ (zeros of
$\sin(t/2)$). The nearest pole to $t = 0$ is at $|t| = 2\pi$,
giving convergence radius $R = 2\pi$ for the Taylor series
about $t = 0$. The Borel transform replaces $t^{2g}$ by
$u^{2g}/(2g)!$, which grows as $|B_{2g}|/(2g)! \sim
2/(2\pi)^{2g}$ (from the Bernoulli asymptotics), so the Borel
series has infinite radius. Since $\hat{Z}$ is entire, the
Laplace integral $\int_0^\infty e^{-u/t}\hat{Z}(u)\,du/t$
converges for all $t > 0$ and reproduces $Z^{\mathrm{sh}}(t)$
term by term. No Stokes phenomenon occurs.
\end{proof}

\subsection{The shadow--Siegel gap}
\label{subsec:shadow-siegel-gap}

The shadow amplitudes $F_g$ and the Igusa cusp form
$\Phi_{10}$ are related but categorically distinct objects.
Four independent obstructions prevent the shadow tower from
reconstructing $\Phi_{10}$.

\begin{theorem}[Shadow--Siegel gap; \ClaimStatusProvedHere]
\label{thm:shadow-siegel-gap}
\index{shadow--Siegel gap|textbf}
\index{Igusa cusp form!shadow gap}
The shadow obstruction tower of $K3 \times E$ does not
produce the Igusa cusp form. The gap consists of four
independent obstructions:
\begin{enumerate}[label=\textup{(O\arabic*)}]
\item \textbf{Categorical.}
$F_g$ is a rational number \textup{(}a tautological
intersection number on $\overline{\mathcal{M}}_g$\textup{)};
$\Phi_{10}(\Omega)$ is a holomorphic function on
$\mathfrak{H}_2$ \textup{(}a Siegel cusp form\textup{)}.
The shadow captures the topological skeleton of the
genus-$g$ amplitude, not the full analytic partition function.
\item \textbf{Modular characteristic.}
The shadow tower uses $\kappa_{\mathrm{ch}} = 3$ (the modular
characteristic of the single-copy chiral algebra
$\Omega^{\mathrm{ch}}(K3 \times E)$, equal to
$\dim_\bC = 3$). The Igusa cusp form has weight
$\operatorname{wt}(\Phi_{10}) = 10 = 2\kappa_{\mathrm{BKM}}$
with $\kappa_{\mathrm{BKM}} = c_0(0)/2 = 10/2 = 5$
(Borcherds weight theorem; Proposition~\textup{\ref{prop:bkm-weight-universal}}). The ratio
$\kappa_{\mathrm{BKM}}/\kappa_{\mathrm{ch}} = 5/3$ reflects
the passage from single-copy to symmetric-product via the
DMVV formula~\textup{\eqref{eq:dvv}}, and is specific to
$K3 \times E$ \textup{(}it is NOT a universal constant:
for the quintic, $\chi = -200$ gives a different ratio\textup{)}.
\item \textbf{Second quantization.}
The shadow tower is first-quantized: $F_g$ is the genus-$g$
amplitude of a \emph{single} copy of $\cA$. The partition
function $1/\Phi_{10}$ is second-quantized: it sums over
\emph{all} symmetric products $\operatorname{Sym}^N(K3)$
via the DMVV formula. The Borcherds lift,
the DMVV infinite product, and the symmetric-product
orbifold construction are second-quantized operations
not captured by the shadow tower alone.
\item \textbf{Schottky at $g \geq 4$.}
The Torelli map $j\colon \mathcal{M}_g \to \mathcal{A}_g$
has image of codimension $(g-2)(g-3)/2$ in~$\mathcal{A}_g$
for $g \geq 4$. Explicitly:
\begin{center}
\small
\renewcommand{\arraystretch}{1.1}
\begin{tabular}{cccccccc}
$g$ & $1$ & $2$ & $3$ & $4$ & $5$ & $6$ & $7$ \\
\hline
$\dim \mathcal{M}_g$ & $1$ & $3$ & $6$ & $9$ & $12$ & $15$ & $18$ \\
$\dim \mathcal{A}_g$ & $1$ & $3$ & $6$ & $10$ & $15$ & $21$ & $28$ \\
$\mathrm{codim}$ & $0$ & $0$ & $0$ & $1$ & $3$ & $6$ & $10$
\end{tabular}
\end{center}
For $g \leq 3$, the Torelli map is an isomorphism
($g = 1$), birational ($g = 2$), or injective with dense
image ($g = 3$), so the shadow tower sees essentially all
Siegel modular form data. At $g = 4$, the Schottky--Igusa
form of weight~$8$ vanishes precisely on the Jacobian locus
$\mathcal{J}_4 = j(\mathcal{M}_4)$: this is information
accessible to Siegel forms but invisible to the shadow tower.
For $g \to \infty$, the ratio
$\mathrm{codim}/\dim\mathcal{A}_g \to 1$: the Jacobian
locus becomes asymptotically negligible.
\end{enumerate}
\end{theorem}

\begin{proof}
(O1) is immediate from the definitions: $F_g \in \bQ$ while
$\Phi_{10}$ is a holomorphic function of three complex variables.
(O2) follows from the two independent computations
$\kappa_{\mathrm{ch}} = \dim_\bC = 3$
(Proposition~\ref{prop:shadow-k3e}) and
$\operatorname{wt}(\Phi_{10}) = 10$
(Definition~\ref{def:igusa-cusp-form}), together with the
identification $\kappa_{\mathrm{BKM}} = c_0(0)/2 = 10/2 = 5$
(Borcherds weight theorem).
(O3) is the content of the DMVV formula: the
second equality in~\eqref{eq:dvv} expresses $1/\Phi_{10}$
as a sum over symmetric products, an operation external to
single-copy chiral algebra theory.
(O4): the codimension formula
$\mathrm{codim}(\mathcal{J}_g, \mathcal{A}_g)
= g(g+1)/2 - (3g-3) = (g-2)(g-3)/2$ for $g \geq 2$ is
classical (Riemann, Schottky).
\end{proof}

\begin{remark}[The genus-$2$ Torelli bridge]
\label{rem:torelli-bridge}
\index{Torelli map!genus-2 bridge}
At genus~$2$, the Torelli map
$j\colon \mathcal{M}_2 \to \mathcal{A}_2$ is birational:
its complement $\mathcal{A}_2 \setminus j(\mathcal{M}_2)$
is the product locus (abelian surfaces isomorphic to a
product of two elliptic curves). This means $\Phi_{10}$
restricts to a well-defined form on $\mathcal{M}_2$
(vanishing on the product locus at the boundary of the
Satake compactification), and the genus-$2$ partition
function $Z_2(\Omega)$ related to $1/\Phi_{10}$ has a
shadow amplitude
\[
F_2 = \int_{\overline{\mathcal{M}}_2} \kappa_{\mathrm{ch}} \cdot
\lambda_2^{\mathrm{FP}} = \frac{7}{1920}
\]
as its topological residue. The shadow $F_2$ is what
remains of $1/\Phi_{10}$ after integrating out all modular
dependence: a tautological intersection number in the
Chow ring of $\overline{\mathcal{M}}_2$.
\end{remark}

\begin{conjecture}[Full genus-$2$ partition function of $K3 \times E$;
\ClaimStatusConjectured]
\label{conj:genus2-k3e-full}
\index{genus-2 partition function!K3 x E (full)}
\index{shadow correction!genus-2}
The full genus-$2$ chiral partition function of $K3 \times E$
decomposes as
\begin{equation}\label{eq:z2-k3e-full}
Z_2^{K3 \times E}(\Omega)
= Z_2^{\mathrm{Heis},\, c=24}(\Omega) \cdot R_{\mathrm{shadow}}^{K3}(\Omega),
\end{equation}
where the Heisenberg base is $Z_2^{\mathrm{Heis},\, c=24}
= \Delta_5(\Omega)^{-2} = 1/\Phi_{10}(\Omega)$
(Siegel modular of weight $-12$ on $\mathrm{Sp}(4,\bZ)$),
and the shadow correction factor is
\begin{equation}\label{eq:r-shadow-k3}
R_{\mathrm{shadow}}^{K3}(\Omega) = 1 + \sum_{k \geq 3} S_k^{K3} \cdot C_k(\Omega).
\end{equation}
Here $S_k^{K3}$ are the shadow tower coefficients of the $K3$ sigma
model VOA \textup{(}the small $\cN = 4$ SCA at $c = 6$, $k_R = 1$;
$\kappa_{\mathrm{ch}} = 2$, $\alpha = 2$, $S_4 = 5/156$, class~$M$\textup{)},
and $C_k(\Omega)$ are genus-$2$ modular cocycles arising from integrating
the $k$-point OPE kernel over~$\Sigma_2$.

The shadow correction $R_{\mathrm{shadow}}^{K3}$ is \emph{not} a Siegel
modular form: it transforms as a mock Siegel modular form of weight~$0$,
with the mock-modular source being the quasi-modular Eisenstein series $E_2$
in the leading cocycle~$C_3$. The shadow series is Gevrey-$1$ divergent
\textup{(}class~$M$: Borel summable but not convergent\textup{)}.

The existence of $\cA_{K3 \times E}$ is established by
CY-A$_3$ \textup{(}proved in the infinity-categorical framework;
Theorem~\ref{thm:derived-framing-obstruction}\textup{)}.
The conjecture's remaining content is the factorisation
structure~\eqref{eq:z2-k3e-full} and the identification of the
shadow correction factor. The shadow tower data of the
$K3$ sigma model is proved at $d = 2$
\textup{(}CY-A$_2$, Theorem~\ref{thm:cy-a-d2}\textup{)}.
The \textup{74} tests in
\texttt{genus2\_k3e\_full.py} verify the factorisation, cocycle
structure, $\mathrm{Sp}(4,\bZ)$ weights, and Borel resummation.
\end{conjecture}

\begin{remark}[Shadow--Siegel bridge at genus~$2$]
\label{rem:shadow-siegel-bridge-g2}
\index{shadow--Siegel gap!genus-2 bridge}
The identity $Z_2^{\mathrm{Heis},\, c=24} = 1/\Phi_{10}$ says that
the Heisenberg base \emph{already is} the DMVV answer.
The shadow correction $R_{\mathrm{shadow}}^{K3}$ therefore represents
first-quantised $A_\infty$ corrections \emph{beyond} the DMVV formula:
the non-perturbative, mock-modular sector invisible to the symmetric
product.

At $z = 0$ \textup{(}separating degeneration $\Sigma_2 \to
\Sigma_1 \cup \Sigma_1$\textup{)}, all cocycles $C_k$ vanish and
$R_{\mathrm{shadow}} = 1$, recovering the DMVV answer.
The departure from $1/\Phi_{10}$ is controlled by the propagator
$|z|^2/\bigl(\operatorname{Im}(\tau)\operatorname{Im}(\sigma)\bigr)$
between the two handles.

The $E_3$ bar cohomology at genus~$2$ for $K3 \times E$
\textup{(}class~$M$\textup{)} has
$\dim H^*(B^{E_3}(\cA)) = 6^2 = 36$
\textup{(}cf.\ $8^2 = 64$ for classes~$L, C$\textup{)},
with Poincar\'e polynomial $(3t(1+t))^2 = 9t^2 + 18t^3 + 9t^4$.
\end{remark}

\subsection{Euler characteristic and Donaldson--Thomas structure}
\label{subsec:euler-dt}

\begin{proposition}[$\chi(K3 \times E) = 0$ and its consequences;
\ClaimStatusProvedElsewhere]
\label{prop:chi-zero}
\index{Euler characteristic!K3 x E}
\index{Donaldson--Thomas!K3 x E}
The Euler characteristic of $K3 \times E$ vanishes:
\[
\chi(K3 \times E) = \chi(K3) \cdot \chi(E)
= 24 \cdot 0 = 0.
\]
This has three structural consequences for the
enumerative geometry:
\begin{enumerate}[label=\textup{(\roman*)}]
\item \emph{MacMahon trivialization.}
The degree-$0$ DT partition function of a Calabi--Yau
threefold $X$ is $Z_{\mathrm{DT},0}(X) = M(-q)^{\chi(X)}$,
where $M(q) = \prod_{n \geq 1}(1-q^n)^{-n}$ is the
MacMahon function counting plane partitions. For
$\chi = 0$: $M(-q)^0 = 1$, so
$Z_{\mathrm{DT},0}(K3 \times E) = 1$. There is no
degree-$0$ contribution from plane-partition-type
combinatorics.
\item \emph{DT/PT correspondence.}
The vanishing $\chi = 0$ implies that the
Donaldson--Thomas and Pandharipande--Thomas partition
functions agree without a MacMahon prefactor correction:
\[
Z_{\mathrm{DT}}(K3 \times E)
= Z_{\mathrm{PT}}(K3 \times E).
\]
This is the ``clean'' case of the DT/PT wall-crossing
formula: the wall-crossing factor $M(-q)^{\chi}$ is trivial.
\item \emph{BCOV anomaly coefficient.}
The BCOV holomorphic anomaly coefficient
$\kappa_{\mathrm{BCOV}} = \chi(X)/24$ vanishes for
$K3 \times E$. The anomaly equation
$\bar{\partial}_{\bar{i}} F_g = \frac{1}{2}
C_{\bar{i}}^{\,jk}(D_j D_k F_{g-1} + \sum_{r=1}^{g-1}
D_j F_r\,D_k F_{g-r})$ is
therefore unobstructed at the topological level,
consistent with the shadow tower being controlled by the
single parameter $\kappa_{\mathrm{ch}} = 3$ without BCOV corrections.
\end{enumerate}
\end{proposition}

\begin{remark}[The four $\kappa_{\mathrm{ch}}$ invariants]
\label{rem:four-kappas}
\index{modular characteristic!four variants for K3 x E}
The preceding analysis exhibits four distinct invariants
that share the name ``$\kappa_{\mathrm{ch}}$'' in different parts of the
literature:
\begin{center}
\small
\renewcommand{\arraystretch}{1.2}
\begin{tabular}{lll}
\textbf{Symbol} & \textbf{Value} & \textbf{Source} \\
\hline
$\kappa_{\mathrm{ch}}(\cA)$ & $3$ & Modular characteristic (Vol~I, Theorem~D) \\
$\kappa_{\mathrm{BCOV}}$ & $0$ & $\chi(K3 \times E)/24$ \\
$\kappa_{\mathrm{MacMahon}}$ & $0$ & Exponent of $M(-q)^{\chi}$ \\
$\kappa_{\mathrm{BKM}}$ & $5$ & $\operatorname{wt}(\Phi_{10})/2 = c_0(0)/2$
\end{tabular}
\end{center}
These coincide for the Virasoro algebra
($\kappa_{\mathrm{ch}} = c/2$ in all r\^oles) but diverge for
$K3 \times E$. The modular characteristic $\kappa_{\mathrm{ch}}(\cA) = 3$
is the intrinsic invariant of the chiral algebra; the others
are geometric proxies valid in restricted contexts.
\end{remark}

\begin{remark}[Shadow depth classification]
\label{rem:bkm-shadow-depth}
\index{shadow depth!K3 x E}
The chiral algebra $\Omega^{\mathrm{ch}}(K3 \times E)$
has central charge $c = 6$, multiple generators of
varying conformal weight, and an infinite tower of OPE
coefficients from the $K3$ sigma model. By the
shadow depth classification
(Vol~I, Definition~\ref{def:shadow-depth-classification}),
$K3 \times E$ belongs to class~M (mixed, shadow depth
$r_{\max} = \infty$): the Borcherds product
formula~\eqref{eq:borcherds-lift} has infinitely many
distinct exponents $c_0(D)$, reflecting an infinite
sequence of independent shadow invariants. This is
consistent with the operadic complexity principle:
the sigma model on a Calabi--Yau target, being a fully
interacting CFT, has non-formal $\Eone$-chiral coassociative
structure at all degrees.
\end{remark}

\begin{remark}[Genus stratification summary]
\label{rem:genus-stratification}
\index{genus stratification!K3 x E}
The bridge between the shadow tower and Siegel modular
forms degrades with genus:
\begin{center}
\small
\renewcommand{\arraystretch}{1.2}
\begin{tabular}{cll}
\textbf{Genus} & \textbf{Torelli} & \textbf{Shadow--Siegel status} \\
\hline
$1$ & Isomorphism & $F_1 = 1/8$ matches $\eta^{-6}$ exactly \\
$2$ & Birational & $F_2 = 7/1920$ is topological residue of $1/\Phi_{10}$ \\
$3$ & Injective, dense & $F_3$ captures most Siegel data \\
$\geq 4$ & Schottky obstructed & Shadow sees codim-$\frac{(g{-}2)(g{-}3)}{2}$ slice
\end{tabular}
\end{center}
The genus-$2$ case is the first one where the Torelli map is
birational and $\Phi_{10}$ (the unique genus-$2$
cusp form of weight~$10$) enters directly. The
four obstructions of
Theorem~\ref{thm:shadow-siegel-gap} quantify exactly
what the shadow tower captures (the topological skeleton
$F_g = \kappa_{\mathrm{ch}} \cdot \lambda_g^{\mathrm{FP}}$, all-weight with cross-channel correction $\delta F_g^{\mathrm{cross}}$) and what it
misses (the analytic, second-quantized, and transverse-Siegel
content).
\end{remark}



